\documentclass{report}

\input{~/dev/latex/template/preamble.tex}
\input{~/dev/latex/template/macros.tex}

\title{\Huge{}}
\author{\huge{Nathan Warner}}
\date{\huge{}}
\fancyhf{}
\rhead{}
\fancyhead[R]{\itshape Warner} % Left header: Section name
\fancyhead[L]{\itshape\leftmark}  % Right header: Page number
\cfoot{\thepage}
\renewcommand{\headrulewidth}{0pt} % Optional: Removes the header line
%\pagestyle{fancy}
%\fancyhf{}
%\lhead{Warner \thepage}
%\rhead{}
% \lhead{\leftmark}
%\cfoot{\thepage}
%\setborder
% \usepackage[default]{sourcecodepro}
% \usepackage[T1]{fontenc}

% Change the title
\hypersetup{
    pdftitle={Embedded C++}
}

\begin{document}
    % \maketitle
        \begin{titlepage}
       \begin{center}
           \vspace*{1cm}
    
           \textbf{Embedded C++}
    
           \vspace{0.5cm}
            
                
           \vspace{1.5cm}
    
           \textbf{Nathan Warner}
    
           \vfill
                
                
           \vspace{0.8cm}
         
           \includegraphics[width=0.4\textwidth]{~/niu/seal.png}
                
           Computer Science \\
           Northern Illinois University\\
           United States\\
           
                
       \end{center}
    \end{titlepage}
    \tableofcontents
    \pagebreak 
    \unsect{General Cpp}
    \begin{itemize}
        \item \textbf{Explicit}: n C++, the explicit keyword is a specifier used to prevent the compiler from performing implicit type conversions and copy-initialization. It acts as a safety guardrail, forcing developers to explicitly state their intent when converting data types, which helps eliminate a massive category of accidental, hard-to-find bugs.        
            \bigbreak \noindent 
            By default, if a class has a constructor that can be called with a single argument, the C++ compiler treats it as a converting constructor. This means the compiler will silently convert the argument's type into your class type whenever it thinks it needs to
            \bigbreak \noindent 
            \begin{cppcode}
                class Wallet {
                public:
                    int dollars;
                    // Converting constructor (No explicit keyword)
                    Wallet(int amount) : dollars(amount) {} 
                };

                void printWalletAmount(const Wallet& w) {
                    std::cout << "Wallet has amount: " << w.dollars << std::endl;
                }

                int main() {
                    // Implicit conversion happens here!
                    // The compiler automatically converts the integer '42' into a Wallet object.
                    printWalletAmount(42); 

                    // This also works via implicit copy-initialization
                    Wallet myWallet = 100; 
                }
            \end{cppcode}
            \bigbreak \noindent 
            If we don't want this implicit conversion, we use \texttt{explicit}.
            \bigbreak \noindent 
            \begin{cppcode}
                class Wallet {
                public:
                    int dollars;

                    // Marked as explicit
                    explicit Wallet(int amount) : dollars(amount) {} 
                };
            \end{cppcode}
            \bigbreak \noindent 
            Now we must either explicitly call the constructor or explicitly cast into a Wallet type.
        \item \textbf{Explicit User-Defined Conversion Operators (C++11)}: Starting with C++11, you can also apply explicit to conversion operators. This prevents a class from automatically downgrading or converting itself into a primitive type when passed around
            \bigbreak \noindent 
            \begin{cppcode}
                class SmartPointer {
                public:
                    bool isValid() const { return true; }

                    // Marked as explicit conversion to bool
                    explicit operator bool() const {
                        return isValid();
                    }
                };

                int main() {
                    SmartPointer ptr;

                    // Okay, loops and conditions are explicit contexts
                    if (ptr) { ... }

                    // Error
                    int x = ptr + 5;
                }
            \end{cppcode}
            \bigbreak \noindent 
            Prior to C++11, explicit was only meaningful for single-argument constructors. With the introduction of braced init-lists ({}), constructors with multiple arguments can also participate in implicit conversions, meaning they should often be marked explicit too
            \bigbreak \noindent 
            explicit by default unless you have an intentional architectural reason to allow implicit conversions (e.g., a custom String class allowing literal const char* conversions). 
            \bigbreak \noindent 
            If you separate your class implementation, the explicit keyword goes only on the declaration inside the class header file, not on the definition
        \item \textbf{NRVO (Named return value optimization)}: NVRO (more commonly abbreviated as NRVO, or Named Return Value Optimization) is a C++ compiler optimization that constructs a named local variable directly in the caller's memory space to avoid copying or moving
            \bigbreak \noindent 
            Without optimization, a local object is created in the function, copied (or moved) to the return location, and then the local object is destroyed.  The compiler builds the named local object directly at the destination address provided by the caller
            \bigbreak \noindent 
            Zero copies and zero moves happen, even if your object has a complex copy constructor. Unlike unnamed return value optimization (RVO) which became mandatory in C++17, named NVRO remains optional for compilers, though all major compilers (GCC, Clang, MSVC) support it
        \item \textbf{std::move use or don't use}
            \begin{enumerate}
                \item 
                    \bigbreak \noindent 
                    \begin{cppcode}
                        class User {
                        public:
                            explicit User(std::string name)
                                : name_{ /* name or std::move(name)? */ } { ... }

                        private:
                            std::string name_;
                        };
                    \end{cppcode}
                    \bigbreak \noindent 
                    You should use std::move(name) here:
                    \begin{itemize}
                        \item If the caller passes a temporary string (rvalue), it is moved into the parameter, and then moved into the member. Total cost: two cheap moves, zero deep copies.
                        \item If the caller passes an existing string (lvalue), it is copied into the parameter, and then moved into the member. Total cost: one copy, one cheap move.
                    \end{itemize}
                    Pass-by-Value + std::move: Passing by value and then moving it into the member variable is the standard modern C++ idiom (frequently called "sink parameters"). It is highly efficient because:
                \item 
                    \bigbreak \noindent 
                    \begin{cppcode}
                        std::string create_message() {
                            std::string message = "Operation completed";
                            return /* message or std::move(message)? */;
                        }
                    \end{cppcode}
                    \bigbreak \noindent 
                    No std::move
                    \bigbreak \noindent 
                    Because message is a local object being returned by value, the compiler can apply named return value optimization (NRVO), constructing it directly in the caller’s destination.
                    \bigbreak \noindent 
                    If NRVO is not applied, the language permits message to be moved automatically. Writing std::move(message) explicitly can prevent NRVO, potentially making the program perform a move that could have been eliminated entirely.
                \item 
                    \bigbreak \noindent 
                    \begin{cppcode}
                        void store(std::vector<std::unique_ptr<int>>& values) {
                            auto number = std::make_unique<int>(42);

                            values.push_back(
                                /* number or std::move(number)? */
                            );
                        }
                    \end{cppcode}
                    \bigbreak \noindent 
                    It should use 
                    \bigbreak \noindent 
                    \begin{cppcode}
                    values.push_back(std::move(number));
                    \end{cppcode}
                    \bigbreak \noindent 
                    std::unique\_ptr cannot be copied because it represents exclusive ownership. std::move(number) allows push\_back to invoke the move constructor, transferring ownership into the vector.
                \item 
                    \bigbreak \noindent 
                    \begin{cppcode}
                        void process(const std::string& source) {
                            std::string copy{
                                /* source or std::move(source)? */
                            };
                        }
                    \end{cppcode}
                    \bigbreak \noindent 
                    No std::move
                    \bigbreak \noindent 
                    std::move(source) would produce a const std::string\&\&. The usual move constructor expects a non-const rvalue reference:
                    \bigbreak \noindent 
                    Moving generally modifies the source, so it cannot accept a const object. Consequently, std::move(source) would still invoke the copy constructor here.
                \item 
                    \bigbreak \noindent 
                    \begin{cppcode}
                        class Buffer {
                        public:
                            Buffer& operator=(Buffer&& other) {
                                if (this != &other) {
                                    release();

                                    data_ = /* other.data_ or std::move(other.data_)? */;
                                    size_ = other.size_;

                                    other.data_ = nullptr;
                                    other.size_ = 0;
                                }

                                return *this;
                            }

                        private:
                            char* data_ = nullptr;
                            std::size_t size_ = 0;

                            void release();
                    };
                    \end{cppcode}
                    \bigbreak \noindent 
                    No std::move.
                    \bigbreak \noindent 
                    Because data\_ is a raw pointer, std::move(other.data\_) would have no useful effect. Moving and copying a raw pointer both copy its address.
                    \bigbreak \noindent 
                    If the member were a std::unique\_ptr<char[]>, then std::move would be necessary:

            \end{enumerate}
        \item \textbf{noexcept constructors}: In C++, you should always prioritize marking move constructors as noexcept, followed by any other constructor that can structurally guarantee it will never throw an exception
            \bigbreak \noindent 
            Failing to mark eligible constructors with noexcept prevents the compiler from optimizing your code and can cause standard library containers to severely underperform
            \bigbreak \noindent 
            You should almost always mark your move constructors as noexcept
            \bigbreak \noindent 
            Standard library containers like std::vector guarantee strong exception safety. When a vector grows and needs to reallocate its internal memory, it must move the existing elements to the new storage
            \bigbreak \noindent 
            If a move constructor is not marked noexcept, the vector cannot risk using it. If an exception throws mid-move, the original data is altered and corrupted. To stay safe, std::vector will fall back to copying your elements, completely defeating the performance benefits of move semantics
        \item \textbf{std::string\_view}: std::string\_view is a lightweight, non-owning view of character data. It usually contains only a pointer and a length. A function accepting it can efficiently read from several sources without allocating a new string:
            \bigbreak \noindent 
            \begin{cppcode}
                void display(std::string_view text);

                display("literal");
                display(std::string{"temporary"});
                display(existing_string);
            \end{cppcode}
            \bigbreak \noindent 
            The risk is that std::string\_view does not own or extend the lifetime of the characters. Storing a view after its source is destroyed or modified can leave it dangling:
        \item \textbf{Object lifetime model}: In C++, an object's lifetime is a strict runtime property defined by the abstract machine. It dictates the exact boundaries within which an object is considered "alive" and valid to access. Accessing an object outside its lifetime results in undefined behavior
            \bigbreak \noindent 
            An object passes through five distinct phases in its complete lifecycle
            \begin{enumerate}
                \item \textbf{Storage Allocation:} Memory with the proper size and alignment is obtained.
                \item \textbf{Initialization (Lifetime Begins):} The constructor completes.
                \item \textbf{Active Use:} The program reads, modifies, or calls methods on the object.
                \item \textbf{Destruction (Lifetime Ends):} The destructor starts running.
                \item \textbf{Storage Deallocation:} The raw memory is released or reused
            \end{enumerate}
        \item \textbf{When does lifetime begin and end}: According to the C++ standard, the precise bounds depend on whether the object has a "trivial" or "non-trivial" initialization/destructor:
            \begin{center}
                \begin{tabularx}{\textwidth}{@{}XXX@{}}
                    \toprule
                    \textbf{Object Type}& \textbf{Lifetime Begins When...}& \textbf{Lifetime Ends When...} \\
                    \midrule
                    Trivial Types(e.g., int, float, trivial structs) & Storage with proper size and alignment is obtained. & The storage is reused, released, or overwritten. \\[5ex]
                    Class Types(with non-trivial constructor/destructor) & Initialization is complete (the constructor finishes execution). & The destructor call starts. \\
                    \bottomrule
                \end{tabularx}
            \end{center}
            An object’s lifetime generally begins when suitably aligned storage has been obtained and its initialization is complete. Merely allocating raw bytes does not necessarily create an object in them; a non-implicit-lifetime class may require construction, often through placement new or std::construct\_at.
            \bigbreak \noindent 
            For a class object, its lifetime ends when its destructor begins, not when it finishes. An object’s lifetime can also end when its storage is released or reused by another object.
        \item \textbf{Construction / destruction order}: Members are initialized in the order they are declared inside the class: first, then second. Their order in the member-initializer list does not control initialization.
            \bigbreak \noindent 
            More completely, construction proceeds as follows:
            \begin{itemize}
                \item Virtual base classes
                \item Direct base classes
                \item Data members in declaration order
                \item Constructor body
            \end{itemize}
        \item \textbf{Init list vs constructor body assignment}: Members should generally be initialized in the constructor’s member-initializer list instead of assigned inside the constructor body?
            \bigbreak \noindent 
            The initializer list constructs each member directly with its intended value:
            \bigbreak \noindent 
            \begin{cppcode}
            Widget::Widget(int n) : value{n} {}
            \end{cppcode}
            \bigbreak \noindent 
            Assignment in the constructor body first initializes the member, often by default construction, and then assigns another value:
            \bigbreak \noindent 
            \begin{cppcode}
                Widget::Widget(int n) {
                    value = n;
                }
            \end{cppcode}
            \bigbreak \noindent 
            That can perform unnecessary work and does not work for members that cannot be assigned after construction. Reference members and const members must be initialized in the initializer list. A member with no accessible default constructor must also be explicitly initialized there. Base-class constructors are selected there as well.
        \item \textbf{Mutable keyword}: The mutable keyword in C++ allows a class data member to be modified even if the containing object is declared as const. It serves as a tool to preserve logical constness, meaning the object appears completely unchanged from the outside, but its internal, non-observable state can still be updated
            \bigbreak \noindent 
            C++ uses mutable in two distinct scenarios: class data members and lambda expressions
            \begin{enumerate}
                \item \textbf{Class data members}: By default, a const member function cannot modify any data members of its class. However, marking a variable as mutable overrides this restriction. 
                    \bigbreak \noindent 
                    Can be only applied to non-static, non-const, and non-reference data members
                    \bigbreak \noindent 
                    \begin{cppcode}
                    class MathWizard {
                    private:
                        double input;
                        mutable double cachedResult = 0.0;
                        mutable bool isCached = false; // Tracks if cache is valid

                    public:
                        MathWizard(double val) : input(val) {}

                        // Declared const because it doesn't change the user-facing state
                        double getExpensiveCalculation() const {
                            if (!isCached) {
                                // Modifying mutable variables inside a const function
                                cachedResult = input * 42.0; // Imagine a heavy math formula here
                                isCached = true;
                                std::cout << "[Computing fresh result]\n";
                            } else {
                                std::cout << "[Returning cached result]\n";
                            }
                            return cachedResult;
                        }
                    };

                    int main() {
                        const MathWizard wizard(10.0); // Object is completely const

                        std::cout << wizard.getExpensiveCalculation() << "\n"; // Computes
                        std::cout << wizard.getExpensiveCalculation() << "\n"; // Uses cache
                    }
                    \end{cppcode}
                    \bigbreak \noindent 
                    Mutable provides no thread safety. If multiple threads increment hits without synchronization, there is a data race and therefore undefined behavior. Thread safety would require something such as:
                \item \textbf{Lambda expressions}: By default, variables captured by a lambda by-value are treated as const inside the lambda body. You cannot modify them unless you explicitly mark the lambda as mutable
                    \bigbreak \noindent 
                    By-value lambda captures are const by default because the lambda’s capture block generates a compiler-defined class, and its operator() method is marked const.
                    \bigbreak \noindent 
                    When you write a lambda, the C++ compiler transforms it behind the scenes into a standard struct or class called a closure object. The variables you capture by value become member variables of this generated object.
                    \bigbreak \noindent 
                    \begin{cppcode}
                        int x = 10;
                        auto my_lambda = [x]() { /* code */ };
                    \end{cppcode}
                    \bigbreak \noindent 
                    The compiler translates that into a unique class that looks roughly like this:
                    \bigbreak \noindent 
                    \begin{cppcode}
                        // What the compiler creates for you:
                        class AnonymousLambdaClosure {
                        private:
                            int x; // The captured variable becomes a data member

                            public:
                            AnonymousLambdaClosure(int ext_x) : x(ext_x) {}

                            // CRITICAL: The function call operator is automatically marked const!
                            void operator()() const { 
                                // Inside here, 'x' is treated as const because the method is const
                            }
                        };

                        AnonymousLambdaClosure my_lambda(x);
                    \end{cppcode}
                    \bigbreak \noindent 
                    Because operator() is a const member function, it is illegal to modify any of the class's internal member variables (like x) inside the function body.
                    \bigbreak \noindent 
                    In functional programming, a lambda is conceptually treated as a pure, stateless function execution block rather than a stateful object. Making the execution operator const aligns lambdas with this philosophy, ensuring that calling a lambda multiple times yields predictable behavior unless you explicitly opt into statefulness using mutable.
            \end{enumerate}
        \item \textbf{Rule of 3, 5, and 0}: The idea is that if you define any special member function for your class you should probably define all of them.
            \begin{itemize}
                \item \textbf{Rule of Three:} destructor, copy constructor, and copy-assignment operator.
                \item \textbf{Rule of Five:} adds the move constructor and move-assignment operator.
                \item \textbf{Rule of Zero:} delegate ownership to RAII types such as std::vector, std::string, and smart pointers so the compiler-generated special members are sufficient.
            \end{itemize}
            These are design guidelines, not requirements to implement every function. After considering them, some operations may appropriately be = default or = delete.
        \item \textbf{Note about copy assignment}: Copy assignment must safely replace existing allocation without leaking it. It should also handle allocation failure without corrupting the original state and account for self-assignment. Copy-and-swap is one common solution.
        \item \textbf{std::swap and std::copy}:  std::swap exchanges the values of two objects, while std::copy transfers a range of elements from one location to another. Both are core components of the C++ Standard Template Library (<utility> and <algorithm> respectively).
            \bigbreak \noindent 
            std::swap takes two objects by reference and swaps their values. In C++11 and later, it typically uses move semantics to avoid heavy copies.
            \bigbreak \noindent 
            Very fast for primitive types (optimized into registers) and large containers (like std::vector, where it only swaps internal pointers, capacity, and size instead of elements)
            \bigbreak \noindent 
            \begin{cppcode}
                int a = 5;
                int b = 10;
                std::swap(a, b); // a is now 10, b is now 5
            \end{cppcode}
            std::copy copies a range of elements from a source location to a destination location.
            \bigbreak \noindent 
            It takes three iterators: first and last defining the source range, and d\_first pointing to the beginning of the destination range.
            \bigbreak \noindent 
            The destination range must already be allocated and large enough to hold the incoming data (or use back-inserters like std::back\_inserter for containers).
            \bigbreak \noindent 
            \begin{cppcode}
                std::vector<int> src = {1, 2, 3};
                std::vector<int> dest(3); // Must be pre-allocated!
                std::copy(src.begin(), src.end(), dest.begin());
                // dest is now {1, 2, 3}
            \end{cppcode}
            \bigbreak \noindent 
            std::copy is the C++ equivalent of memcpy, but it is much safer and smarter.
        \item \textbf{Copy and swap (copy assignment)}: 
            \bigbreak \noindent 
            \begin{cppcode}
                class Buffer {
                    std::size_t size_ = 0;
                    int* data_ = nullptr;

                    public:
                    Buffer(std::size_t size)
                        : size_{size}, data_{new int[size]} {}

                    Buffer(const Buffer& other)
                        : size_{other.size_}, data_{new int[other.size_]} {
                        std::copy(
                        other.data_,
                        other.data_ + size_,
                        data_
                        );
                    }

                    ~Buffer() {
                        delete[] data_;
                    }
            \end{cppcode}
            \bigbreak \noindent 
            \begin{cppcode}
                friend void swap(Buffer& first, Buffer& second) noexcept {
                    using std::swap;
                    swap(first.size_, second.size_);
                    swap(first.data_, second.data_);
                }

                Buffer& operator=(Buffer other) {
                    swap(*this, other);
                    return *this;
                }
            };
            \end{cppcode}
            \bigbreak \noindent 
            The assignment parameter is deliberately passed by value. For a = b, other is first copy-constructed from b. If allocation fails during that copy, the assignment body never starts and a remains unchanged. Once the copy succeeds:
            \begin{enumerate}
                \item a exchanges resources with other.
                \item a now owns the copied data.
                \item other owns a’s old data.
                \item other is destroyed and releases the old data.
            \end{enumerate}
            Self-assignment also works. a = a creates an independent copy before anything is changed, swaps it into a, and safely destroys the old allocation.
        \item \textbf{Why is marking a move constructor noexcept particularly important when objects of that type are stored in a std::vector? During vector reallocation, why might the vector copy elements instead of moving them?}: During reallocation, std::vector must transfer its elements into new storage. If it moves several elements and a later move throws, the earlier source elements may already have been modified, making rollback difficult.
            \bigbreak \noindent 
            To preserve the strong exception guarantee, the vector generally copies when:
            \begin{itemize}
                \item The move constructor might throw, and
                \item The type is copy-constructible.
            \end{itemize}
            A noexcept move constructor tells the vector that moving cannot fail, allowing it to use the usually cheaper move operation. If the type is move-only, the vector may have no choice but to use a potentially throwing move.
        \item \textbf{Virtual inheritance and the diamond problem}: The diamond problem occurs when a class inherits from two classes that share the same base class.
            \begin{verbatim}
                Device
               /      \
           Scanner   Printer
               \      /
              Copier
            \end{verbatim}
            \bigbreak \noindent 
            In ordinary multiple inheritance, Copier contains two separate Device base-class subobjects: one inherited through Scanner and another through Printer
            \bigbreak \noindent 
            \begin{cppcode}
            class Device {
                public:
                std::string name;
            };

            class Scanner : public Device {};
            class Printer : public Device {};

            class Copier : public Scanner, public Printer {};

            int main() {
                Copier copier;

                copier.name = "Office copier"; // Error: ambiguous
            }
            \end{cppcode}
            \bigbreak \noindent 
            The compiler cannot determine which Device::name you mean:
            \bigbreak \noindent 
            \begin{cppcode}
                copier.Scanner::name = "Scanner side";
                copier.Printer::name = "Printer side";
            \end{cppcode}
            \bigbreak \noindent 
            These are genuinely two different variables because the object conceptually looks like this:
            \bigbreak \noindent 
            \begin{verbatim}
                Copier
                |-- Scanner
                |   |-- Device
                |       |-- name
                |-- Printer
                    |-- Device
                    |-- name
            \end{verbatim}
            \bigbreak \noindent 
            This duplication is often undesirable. A real copier is usually intended to be one device that supports scanning and printing, not two independent devices.
            \bigbreak \noindent 
            Virtual inheritance tells C++ that all relevant inheritance paths should share one base-class subobject.
            \bigbreak \noindent 
            The intermediate classes must inherit the common base using virtual:
            \bigbreak \noindent 
            \begin{cppcode}
                class Scanner : virtual public Device {};
                class Printer : virtual public Device {};

                class Copier : public Scanner, public Printer {};
            \end{cppcode}
            \bigbreak \noindent 
            Now the object conceptually looks like this:
            \begin{verbatim}
                Copier
                |-- Scanner --|
                |-- Printer --|-- shared Device
                |-------------|       |-- name
            \end{verbatim}
            \bigbreak \noindent 
            Consequently, accessing name is no longer ambiguous:
            \bigbreak \noindent 
            \textbf{Note:} Both paths must use virtual inheritance. Virtual inheritance must be specified by every inheritance path that is supposed to share the base
            \bigbreak \noindent 
            One especially important rule is that the most-derived class constructs the virtual base.
            \bigbreak \noindent 
            When a Copier is created, it is Copier, not Scanner or Printer, that constructs the shared Device:
            \bigbreak \noindent 
            Consider 
            \begin{verbatim}
                    A
                    |
                    B
                   / \
                  C   D
                   \ /
                    E
            \end{verbatim}
            \bigbreak \noindent 
            The diamond converges at B, so C and D should inherit virtually from B:
            \bigbreak \noindent 
            \begin{cppcode}
            class A {};

            class B : public A {};

            class C : virtual public B {};
            class D : virtual public B {};

            class E : public C, public D {};
            \end{cppcode}
            \bigbreak \noindent 
            Now an E object contains:
            \begin{itemize}
                \item One shared B subobject
                \item One A subobject contained inside that shared B
            \end{itemize}
            \begin{itemize}
                \item B is a virtual base of C, D, and ultimately E.
                \item A is an ordinary base of B.
                \item E has only one A because it has only one shared B, and that B contains one A.
            \end{itemize}
            Also remember that because B is virtual, E, as the most-derived class, is responsible for constructing it:
            \bigbreak \noindent 
            \begin{cppcode}
                class B : public A {
                    public:
                    explicit B(int value) {}
                };

                class C : virtual public B {
                    public:
                    C() : B(1) {}
                };

                class D : virtual public B {
                    public:
                    D() : B(2) {}
                };

                class E : public C, public D {
                    public:
                    E() : B(3) {} // E constructs the shared virtual B
                };
            \end{cppcode}
            The B(1) and B(2) initializers are ignored while constructing an E; E’s B(3) initializes the single shared B.
            \bigbreak \noindent 
            Both the C and D constructors are still called. Virtual inheritance only changes who constructs the shared B base.
            \bigbreak \noindent 
            When constructing an E object:
            \begin{itemize}
                \item E first constructs its virtual base B.
                    \begin{itemize}
                        \item Before B’s constructor body runs, B constructs its ordinary base A.
                    \end{itemize}
                \item E constructs its first direct base, C.
                \item E constructs its second direct base, D.
                \item E constructs its own data members in declaration order.
                \item The body of E’s constructor runs.
            \end{itemize}
            Destruction occurs in the exact reverse order:
            \bigbreak \noindent 
            \begin{cppcode}
                Destroy E
                Destroy D
                Destroy C
                Destroy B
                Destroy A
            \end{cppcode}
            \bigbreak \noindent 
            B is constructed only once because it is a virtual base shared by C and D.
        \item \textbf{C++ Casting}: 
            \begin{itemize}
                \item \textbf{static\_cast}: Performs compile-time conversions between related types.
                \item \textbf{dynamic\_cast}: Safely converts pointers or references within a class hierarchy at runtime.
                    \bigbreak \noindent 
                    Commonly used when downcasting from a base class pointer to a derived class pointer.
                    \bigbreak \noindent 
                    Requires the base class to have at least one virtual function (polymorphism). Returns nullptr (for pointers) or throws an exception (for references) if the cast fails.
                \item \textbf{const\_cast}: Adds or removes the const or volatile qualifier from a variable.
                    \bigbreak \noindent 
                    Commonly used when passing a const variable to a legacy API that accepts a non-const pointer but does not modify the data.
                    \bigbreak \noindent 
                    Modifying an object that was originally declared as const after casting away its constness results in undefined behavior.
                \item \textbf{reinterpret\_cast}: Reinterprets the underlying bit pattern of an object as a completely different type.
                    \bigbreak \noindent 
                     Low-level programming, such as converting a pointer to an integer (uintptr\_t) or casting raw byte buffers.
            \end{itemize}
            dynamic\_cast is the only C++ cast that performs runtime checking.






    \end{itemize}

    \pagebreak 
    \subsection{Smart pointers}
    \begin{itemize}
        \item \textbf{Why is std::unique\_ptr<T> non-copyable?}: Copying would violate exclusive ownership. A unique pointer represents one owning handle, so copying it would create two owners. 
        \item \textbf{How should ownership normally be transferred from one std::unique\_ptr to another?}: Move the unique pointer.
        \item \textbf{What normally happens when a nonempty std::shared\_ptr is copied?}: The strong reference count increases
        \item \textbf{What is the principal ownership purpose of std::weak\_ptr?}: Observe shared ownership without extending lifetime
        \item \textbf{What does calling lock() on a std::weak\_ptr produce?}: A shared pointer, empty if expired
        \item \textbf{Which operation makes a std::unique\_ptr stop owning its object without deleting that object?}: release() returns the raw pointer and leaves the unique pointer empty without invoking its deleter
        \item \textbf{What deletion operation does std::unique\_ptr<T[]> use by default?}: delete[]
        \item \textbf{What is wrong with constructing std::shared\_ptr<T> second(first.get()) when first already owns the object?}: It can cause the object to be deleted twice
            \bigbreak \noindent 
            The two shared pointers use separate control blocks, so each may independently delete the same raw pointer
        \item \textbf{Why can a cycle made entirely of std::shared\_ptr objects leak memory?}: Every object in the cycle can remain strongly owned
            \bigbreak \noindent 
            Internal strong references can prevent every count in the cycle from reaching zero after external owners disappear.
        \item \textbf{Which change commonly breaks an unintended std::shared\_ptr ownership cycle?}: Replace one non-owning relationship with std::weak\_ptr.
        \item \textbf{Cycle problem}: Suppose two objects own each other with std::shared\_ptr. Neither reference count can reach zero, so neither object is destroyed.
            \bigbreak \noindent 
            \begin{cppcode}
                class Person {
                    public:
                    std::string name;
                    std::shared_ptr<Person> partner;

                    explicit Person(std::string name) : name(std::move(name)) {
                        std::cout << name << " created\n";
                    }

                    ~Person() { std::cout << name << " destroyed\n"; }
                };
            \end{cppcode}
            \bigbreak \noindent 
            \begin{cppcode}
                int main() {
                    auto alice = std::make_shared<Person>("Alice");
                    auto bob   = std::make_shared<Person>("Bob");

                    alice->partner = bob;
                    bob->partner   = alice;
                }
            \end{cppcode}
            \bigbreak \noindent 
            Notice that there are two shared pointers pointing to the same memory, so the reference count for both Person instances is two.
            \bigbreak \noindent 
            When main() ends:
            \begin{enumerate}
                \item The bob smart-pointer variable is destroyed, bobs reference count decreases to one, bob not destroyed because alice->partner still owns it.
                \item The alice smart-pointer variable is destroyed, alice reference count decreases to one, alice not destroyed because bob->partner still owns it.
            \end{enumerate}
            Therefore neither Bob nor Alice have their destructors called. Each object keeps the other alive, so neither reference count reaches zero.
            \bigbreak \noindent 
            With weak\_ptr, the partner references do not increase the strong reference counts. Therefore, destroying bob reduces Bob’s count from one to zero, and destroying alice does the same for Alice. Both Person objects are then correctly destroyed.
        \item \textbf{Unique pointer core api}: 
            \begin{itemize}
                \item \textbf{Creation}:  
                    \bigbreak \noindent 
                    \begin{cppcode}
                        auto p = std::make_unique<MyType>(arg1, arg2);
                        auto array = std::make_unique<int[]>(100);
                    \end{cppcode}
                    \bigbreak \noindent 
                    Prefer std::make\_unique over writing new directly.
                \item \textbf{Access}:
                    \bigbreak \noindent 
                    \begin{cppcode}
                        p->member;
                        *p;
                        p.get();             // Non-owning raw pointer
                        static_cast<bool>(p);
                        if (p) { }
                    \end{cppcode}
                    \bigbreak \noindent 
                    get() does not transfer ownership. The returned pointer must not normally be deleted.
                \item \textbf{Transfer ownership}:
                    \bigbreak \noindent 
                    \begin{cppcode}
                    auto p2 = std::move(p);
                    \end{cppcode}
                    \bigbreak \noindent 
                    After the move, p is normally empty. A unique pointer cannot be copied
                \item \textbf{Replace or destroy}:
                    \bigbreak \noindent 
                    \begin{cppcode}
                        p.reset();            // Delete object and become empty
                        p.reset(new MyType);  // Delete old object and own new object
                        p = nullptr;          // Delete object and become empty
                    \end{cppcode}
                    \bigbreak \noindent 
                    Prefer assigning std::make\_unique when replacing an object:
                    \bigbreak \noindent 
                    \begin{cppcode}
                    p = std::make_unique<MyType>();
                    \end{cppcode}
                \item \textbf{Release ownership}:
                    \bigbreak \noindent 
                    \begin{cppcode}
                    MyType* raw = p.release();
                    \end{cppcode}
                    \bigbreak \noindent 
                    release() makes p empty without deleting the object. The caller becomes responsible for it:
                \item \textbf{Swap}:
                    \bigbreak \noindent 
                    \begin{cppcode}
                        p.swap(p2);
                        std::swap(p, p2);
                    \end{cppcode}
                \item \textbf{Custom deleter}:
                    \bigbreak \noindent 
                    \begin{cppcode}
                        struct FileCloser {
                            void operator()(FILE* file) const {
                                if (file)
                                std::fclose(file);
                            }
                        };

                        std::unique_ptr<FILE, FileCloser> file(
                            std::fopen("data.txt", "r")
                        );
                    \end{cppcode}
                    \bigbreak \noindent 
                    The deleter is part of the unique\_ptr type:
                    \bigbreak \noindent 
                    \begin{cppcode}
                    std::unique_ptr<FILE, FileCloser>
                    \end{cppcode}
            \end{itemize}
        \item \textbf{Shared pointer core api}
            \begin{itemize}
                \item \textbf{Creation}:  
                    \bigbreak \noindent 
                    \begin{cppcode}
                    auto p = std::make_shared<MyType>(arg1, arg2);
                    \end{cppcode}
                    \bigbreak \noindent 
                    You can copy a shared pointer:
                    \bigbreak \noindent 
                    \begin{cppcode}
                    You can copy a shared pointer:
                    \end{cppcode}
                \item \textbf{Access}:
                    \bigbreak \noindent 
                    \begin{cppcode}
                        p->member;
                        *p;
                        p.get();             // Non-owning raw pointer
                        if (p) { }
                    \end{cppcode}
                \item \textbf{Inspect the strong reference count}:
                    \bigbreak \noindent 
                    \begin{cppcode}
                    long count = p.use_count();
                    \end{cppcode}
                \item \textbf{Release one owner}:
                    \bigbreak \noindent 
                    \begin{cppcode}
                        p.reset();
                        p = nullptr;
                    \end{cppcode}
                    \bigbreak \noindent 
                    This decreases the strong reference count. The object is destroyed only when the count reaches zero.
                \item \textbf{Replace the managed object}:
                    \bigbreak \noindent 
                    \begin{cppcode}
                    p = std::make_shared<MyType>();
                    p.reset(new MyType);
                    \end{cppcode}
                \item \textbf{Move ownership}: 
                    \bigbreak \noindent 
                    \begin{cppcode}
                    auto p2 = std::move(p);
                    \end{cppcode}
                    \bigbreak \noindent 
                    Moving does not increase the reference count. Ownership represented by p is transferred to p2.
                \item \textbf{Swap}:
                    \bigbreak \noindent 
                    \begin{cppcode}
                        p.swap(p2);
                        std::swap(p, p2);
                    \end{cppcode}
                \item \textbf{Convert from unique\_ptr}: A unique\_ptr can transfer ownership into a shared\_ptr:
                    \bigbreak \noindent 
                    \begin{cppcode}
                        std::unique_ptr<MyType> unique =
                        std::make_unique<MyType>();

                        std::shared_ptr<MyType> shared = std::move(unique);
                    \end{cppcode}
                    \bigbreak \noindent 
                    The reverse conversion is not generally available because shared ownership cannot safely become exclusive without additional guarantees.
                \item \textbf{Pointer casts}: Use smart-pointer aware casts
                    \bigbreak \noindent 
                    \begin{cppcode}
                        auto derived =
                        std::static_pointer_cast<Derived>(base);

                        auto checked =
                        std::dynamic_pointer_cast<Derived>(base);

                        auto constant =
                        std::const_pointer_cast<const MyType>(p); 
                    \end{cppcode}
                    \bigbreak \noindent 
                    These preserve the shared control block. 
                    \bigbreak \noindent 
                    C++ provides four primary pointer cast functions for std::shared\_ptr via <memory>
                    \begin{itemize}
                        \item \textbf{std::static\_pointer\_cast:} Casts a pointer up or down an inheritance hierarchy at compile time using static\_cast. It performs no runtime checks, so you must ensure the conversion is valid
                        \item \textbf{std::dynamic\_pointer\_cast:} Safely downcasts a polymorphic base class pointer to a derived class pointer at runtime using dynamic\_cast. It returns an empty std::shared\_ptr if the object is not of the target type. 
                        \item \textbf{std::const\_pointer\_cast:} Adds or removes const qualifiers from the managed object type using const\_cast
                        \item \textbf{std::reinterpret\_pointer\_cast:} Reinterprets the underlying pointer type to an unrelated type using reinterpret\_cast. This is low-level and rarely needed
                    \end{itemize}
                \item \textbf{Custom deleter}:
                    \bigbreak \noindent 
                    \begin{cppcode}
                        std::shared_ptr<FILE> file(
                            std::fopen("data.txt", "r"),
                            [](FILE* file) {
                                if (file)
                                std::fclose(file);
                            }
                        );
                    \end{cppcode}
                    \bigbreak \noindent 
                    Unlike unique\_ptr, the deleter is not part of the shared\_ptr variable’s type.
                    \begin{itemize}
                        \item \textbf{std::unique\_ptr<T, Deleter>} requires the deleter type and is part of the class type.
                        \item \textbf{std::shared\_ptr<T>} does not include the deleter in its template parameters; you pass the deleter to the constructor instead
                            \bigbreak \noindent 
                            std::shared\_ptr already uses a heap-allocated control block for reference counts, so storing a type-erased deleter inside that control block adds no extra pointer overhead per smart pointer instance
                    \end{itemize}
            \end{itemize}
        \item \textbf{Weak pointer core api}: Use to observe an object managed by shared\_ptr without keeping it alive.
            \begin{itemize}
                \item \textbf{Creation}: 
                    \bigbreak \noindent 
                    \begin{cppcode}
                        std::shared_ptr<MyType> shared =
                        std::make_shared<MyType>();

                        std::weak_ptr<MyType> weak = shared;
                    \end{cppcode}
                \item \textbf{Safely access the object}:
                    \bigbreak \noindent 
                    \begin{cppcode}
                        if (auto locked = weak.lock()) {
                            locked->do_something();
                        }
                    \end{cppcode}
                    \bigbreak \noindent 
                    lock() returns:
                    \begin{itemize}
                        \item A nonempty shared\_ptr if the object is still alive.
                        \item An empty shared\_ptr if the object has been destroyed.
                    \end{itemize}
                    The returned shared\_ptr keeps the object alive while it is being used.
                \item \textbf{Test for expiration}:
                    \bigbreak \noindent 
                    \begin{cppcode}
                        if (weak.expired()) {
                            // Object has been destroyed
                        }
                    \end{cppcode}
                \item \textbf{Inspect the strong count}:
                    \bigbreak \noindent 
                    \begin{cppcode}
                    long count = weak.use_count();
                    \end{cppcode}
                \item \textbf{Stop observing}:
                    \bigbreak \noindent 
                    \begin{cppcode}
                    weak.reset();
                    \end{cppcode}
                \item \textbf{std::enable\_shared\_from\_this}: Use when an object managed by shared\_ptr needs to produce another shared\_ptr to itself.
                    \bigbreak \noindent 
                    \begin{cppcode}
                        class Widget : public std::enable_shared_from_this<Widget> {
                        public:
                            std::shared_ptr<Widget> get_pointer() {
                                return shared_from_this();
                            }
                        };

                        auto widget = std::make_shared<Widget>();
                        auto another_owner = widget->get_pointer();
                    \end{cppcode}
                    \bigbreak \noindent 
                    Do not construct a new shared\_ptr from this:
                    \bigbreak \noindent 
                    \begin{cppcode}
                    std::shared_ptr<Widget>(this); // Dangerous
                    \end{cppcode}
                    \bigbreak \noindent 
                    That creates a separate control block and can cause double deletion. shared\_from\_this() requires the object to already be managed by an appropriate shared\_ptr. Otherwise, it throws std::bad\_weak\_ptr.
                    \bigbreak \noindent 
                    C++ 17 also provides
                    \bigbreak \noindent 
                    \begin{cppcode}
                    std::weak_ptr<Widget> weak = weak_from_this();
                    \end{cppcode}

            \end{itemize}
        \item \textbf{make\_shared and make\_unique}: You should use std::make\_unique (and make\_shared) (introduced in C++14) over direct initialization with unique\_ptr<T>(new T()) for four primary reasons
            \bigbreak \noindent 
            std::make\_unique ensures atomicity between allocating memory and wrapping it in a smart pointer. Consider this function call:
            \bigbreak \noindent 
            \begin{cppcode}
            foo(std::unique_ptr<A>(new A()), std::unique_ptr<B>(new B()));
            \end{cppcode}
            \bigbreak \noindent 
            Prior to C++17, the compiler could reorder evaluations. It might allocate memory for A, then allocate memory for B, and then run the constructors. If B's constructor threw an exception, the memory allocated for A would leak because its unique\_ptr hadn't been constructed yet
            \bigbreak \noindent 
            \begin{cppcode}
            foo(std::make_unique<A>(), std::make_unique<B>());
            \end{cppcode}
            \bigbreak \noindent 
            With std::make\_unique, the allocation and construction happen in an isolated, atomic step. If an exception occurs, memory is cleaned up instantly.
            \bigbreak \noindent 
            Modern C++ advocates for total encapsulation of raw memory management. Prior to C++14, the rule of thumb was "Never use new unless you are creating a unique\_ptr". By using std::make\_unique, you can enforce a clean, universal guideline across your entire codebase: Never use the new keyword explicitly



    \end{itemize}

    \pagebreak 
    \subsection{Rvalue references and perfect forwarding}
    \begin{itemize}
        \item \textbf{Rvalue references}: Before C++11, passing or returning massive objects (like a std::vector with millions of elements) often forced the compiler to deeply copy all the data
            \bigbreak \noindent 
            With rvalue references, if the compiler knows an object is a temporary (an rvalue), it can trigger a move constructor instead of a copy constructor. Instead of duplicating the heap data, the new object simply "steals" the internal pointers of the temporary object, reducing the operation's computational cost to near-zero
            \bigbreak \noindent 
            \begin{cppcode}
                // A typical move constructor uses an rvalue reference
                Widget(Widget&& other) noexcept {
                    this->data = other.data;  // Steal the pointer
                    other.data = nullptr;     // Clean up the temporary object so its destructor is safe
                }
            \end{cppcode}
            \bigbreak \noindent 
            If you have an lvalue (a named variable) that you are finished with and want to treat like a temporary to trigger a move, you can cast it using std::move. std::move does not actually move anything; it simply casts an lvalue into an rvalue reference
            \bigbreak \noindent 
            \begin{cppcode}
                std::vector<int> v1 = {1, 2, 3};
                std::vector<int> v2 = std::move(v1); // v2 steals v1's data. v1 is now empty.
            \end{cppcode}
            \bigbreak \noindent 
            This is the most common trap when working with rvalue references: \textbf{If an rvalue reference has a name, the expression evaluating it is an lvalue.}
            \bigbreak \noindent 
            \begin{cppcode}
                void process(std::vector<int>&& param) {
                    // 'param' is an rvalue reference type, but the *expression* 'param' is an lvalue!

                    std::vector<int> internal1 = param;             // TRAP: Triggers a COPY constructor!
                    std::vector<int> internal2 = std::move(param);  // CORRECT: Triggers the MOVE constructor.
                }
            \end{cppcode}
        \item \textbf{Perfect forwarding}: When writing generic wrapper functions (like factory functions, std::make\_unique, or thread constructors), you often need to take arguments and pass them exactly as they are to another function.
            \bigbreak \noindent 
            To do this perfectly, you must preserve:
            \begin{itemize}
                \item \textbf{The value category:} If the user passed an rvalue, the target function should get an rvalue (to trigger a move). If they passed an lvalue, it should get an lvalue.
                \item \textbf{Const modifiers:} If it was const, it must stay const
            \end{itemize}
            Before C++11, achieving this required a massive explosion of function overloads to cover every combination of const, non-const, lvalue, and rvalue references.
            \bigbreak \noindent 
            A forwarding reference (originally termed a universal reference by Scott Meyers) looks exactly like an rvalue reference (T\&\&), but it behaves completely differently. It can bind to lvalues (as a normal reference) and rvalues (as an rvalue reference), preserving constness.
            \bigbreak \noindent 
            A reference becomes a forwarding reference only when the two following conditions are met:
            \begin{itemize}
                \item \textbf{Type deduction is involved:} The type must be deduced by the compiler.
                \item \textbf{The syntax is exactly T\&\&:} It cannot have modifiers like const T\&\&
            \end{itemize}
            \bigbreak \noindent 
            \begin{cppcode}
                template<typename T>
                void wrapper(T&& arg) { // Deduced type + T&& syntax = Forwarding Reference!
                    // ...
                }
            \end{cppcode}
            \bigbreak \noindent 
            The following are NOT universal references.
            \bigbreak \noindent 
            \begin{cppcode}
                void foo(int&& arg); // Not a template. This is strictly an rvalue reference.

                template<typename T>
                void bar(const T&& arg); // Has 'const'. This is strictly an rvalue reference.

                template<typename T>
                class Widget {
                    void baz(T&& arg); // 'T' belongs to the class, not the function. 
            };                     // No type deduction occurs when calling baz(), so it's an rvalue reference.
            \end{cppcode}
            \bigbreak \noindent 
            How can T\&\& magically bind to an lvalue? The secret lies in a compiler mechanism called reference collapsing.
            \bigbreak \noindent 
            In C++, you cannot write a reference to a reference directly (e.g., int\& \&\& will fail to compile). However, the compiler generates them behind the scenes during template type deduction. To resolve them, it applies a simple rule: An lvalue reference always wins.
            \bigbreak \noindent 
            \begin{itemize}
                \item \&  + \&  $\to$ \&
                \item \&  + \&\& $\to$ \&
                \item \&\& + \&  $\to$ \&
                \item \&\& + \&\& $\to$ \&\&
            \end{itemize}
            \bigbreak \noindent 
            \begin{cppcode}
                template<typename T>
                void wrapper(T&& arg);

                int x = 10;
                wrapper(x); // 1. You pass an lvalue (x). 
                // 2. Compiler deduces T as 'int&'.
                // 3. Substituting T gives: 'int& && arg'
                // 4. Collapsing rule applies: 'int& arg' (an lvalue reference!)

                wrapper(20); // 1. You pass an rvalue (20).
                // 2. Compiler deduces T as 'int'.
                // 3. Substituting T gives: 'int&& arg' (an rvalue reference!)
            \end{cppcode}
        \item \textbf{std::forward}: As established with rvalue references, once a reference has a name, it evaluates as an lvalue. Inside our wrapper function, the variable arg has a name. Therefore, if we just pass arg to a target function, it will always be treated as an lvalue, defeating the whole purpose!
            \bigbreak \noindent 
            Unlike std::move (which unconditionally turns an object into an rvalue), std::forward<T> conditionalizes its behavior based on how T was deduced:
            \begin{itemize}
                \item If T was deduced as an lvalue reference (int\&), std::forward casts it back to an lvalue reference.
                \item If T was deduced as a normal type (int), std::forward casts it back to an rvalue reference.
            \end{itemize}
            \bigbreak \noindent 
            \begin{cppcode}
                #include <iostream>
                #include <utility>

                void target(int& x) { std::cout << "Lvalue overload\n"; }
                void target(int&& x) { std::cout << "Rvalue overload\n"; }

                template<typename T>
                void perfectWrapper(T&& arg) {
                    // std::forward restores the original value category
                    target(std::forward<T>(arg)); 
                }

                int main() {
                    int a = 5;

                    perfectWrapper(a);  // Prints: Lvalue overload
                    perfectWrapper(10); // Prints: Rvalue overload
                }
            \end{cppcode}
        \item \textbf{Type vs value category}: In C++, every single expression is characterized by two completely independent properties: its type and its value category
            \begin{itemize}
                \item \textbf{Type} answers the question: "What kind of data is this, and what operations can I perform on it?" 
                \item \textbf{Value category}: Value Category answers the questions: "Does this data have a permanent address in memory (identity), and can I safely steal its resources (movability)?
            \end{itemize}
            \bigbreak \noindent 
            Value categories are a property of expressions, not variables or objects. An expression is any piece of code that evaluates to a value (like x, 5 + 3, or foo())
            \bigbreak \noindent 
            Modern C++ (C++11 and later) classifies expressions based on two core criteria:
            \begin{itemize}
                \item \textbf{Identity}: Does the program have a way to determine the exact memory address of the object
                \item \textbf{Movability}: Can the object's resources be safely hijacked/moved because its lifetime is ending? 
            \end{itemize}
        \item \textbf{The 3 Primary Value Categories}: Every expression belongs to exactly one of these three primary buckets: 
            \begin{itemize}
                \item \textbf{lvalue (Left-hand value):} Has identity but cannot be safely moved. These are typically names of variables or things that persist in memory.
                \item \textbf{prvalue (Pure rvalue):} Has no identity but can be moved. These are pure mathematical values or temporary objects that do not exist yet beyond the evaluation.
                \item  \textbf{xvalue (eXpiring value):} Has identity and can be moved. It points to a real place in memory, but the developer has explicitly given permission to strip its resources because it is about to expire.
            \end{itemize}
        \item \textbf{Mixed categories}: The C++ standard groups these primary categories into two broader term
            \begin{itemize}
                \item \textbf{glvalue (Generalized lvalue):} Anything with identity (lvalue + xvalue).
                \item \textbf{rvalue:} Anything that can be moved (prvalue + xvalue)
            \end{itemize}
            The single most confusing concept in C++ is that a variable's type is independent of the value category of the expression containing its name
            \bigbreak \noindent 
            \begin{cppcode}
                void process(int&& arg) {
                    // The TYPE of 'arg' is an rvalue reference (int&&)
                    // The VALUE CATEGORY of the expression 'arg' is an lvalue!
                    int y = arg; 
                }
            \end{cppcode}
        \item \textbf{auto deduction vs CTAD}: While both features automate type resolution in C++, they operate at completely different stages of object creation. auto deduces the type of a variable from its initializer expression, whereas CTAD (Class Template Argument Deduction) deduces a class template’s missing parameters from its constructor arguments.
            \begin{itemize}
                \item \textbf{Auto}: Deduces the complete type of a variable, uses \textbf{template function type deduction} rules
                    \bigbreak \noindent 
                    It strips top-level const, volatile, and references by default, unless you explicitly decorate it (e.g., const auto\&).
                \item \textbf{CTAD}: Deduces missing template arguments for a class. Uses \textbf{constructor signatures} and implicit / explicit deduction guides
                    \bigbreak \noindent 
                    Before C++17, you always had to explicitly state template arguments when instantiating a class, or rely on helper factory functions like std::make\_pair.
                    \bigbreak \noindent 
                    With CTAD, you can omit the template brackets <...> entirely. The compiler looks at the arguments passed to the constructor and figures out what the template parameters should be.
                    \bigbreak \noindent 
                    \begin{cppcode}
                        // Pre-C++17
                        std::pair<int, double> p1(1, 2.3); 
                        auto p2 = std::make_pair(1, 2.3); // Workaround using auto

                        // Post-C++17 (CTAD)
                        std::pair p3(1, 2.3);             // Automatically deduces std::pair<int, double>
                    \end{cppcode}
            \end{itemize}
            You will frequently see auto and CTAD used in the exact same line of code. They do not conflict; they handle two separate responsibilities:
            \bigbreak \noindent 
            \begin{cppcode}
            auto v = std::vector{1, 2, 3};
            \end{cppcode}
            \begin{itemize}
                \item \textbf{CTAD steps in first:} It sees std::vector\{1, 2, 3\}, looks at the integer arguments, and resolves the class type to std::vector<int>.
                \item \textbf{auto steps in second:} It looks at the resolved right-hand side (std::vector<int>) and assigns that exact type to the variable v.
            \end{itemize}
        \item \textbf{auto vs decltype}: auto and decltype both perform type inference at compile time, but they use different deduction rules and serve distinct purposes
            \bigbreak \noindent 
            The core difference is that auto deduces a type based on an initialization value (similar to template argument deduction), whereas decltype yields the exact declared type of an expression or variable without evaluating it
            \begin{itemize}
                \item \textbf{Auto}: Declaring a variable whose type matches its initializer.
                \item \textbf{decltype(expr)}: Inspecting the exact type of a variable or an expression.
                    \bigbreak \noindent 
                    decltype also preserves references and const. The important thing to note is that decltype does not require an initializer, and does not evaluate the expression.
            \end{itemize}
        \item \textbf{Questions}: 
            \begin{enumerate}
                \item \textbf{What kinds of expressions can bind to int\&\&? How does that differ from what can bind to const int\&?}: Both pure rvalues (prvalues) and lvalues cast to an rvalue reference (with std::move) can bind to int\&\&. 
                    \bigbreak \noindent 
                    Both prvalues and lvalues can bind to int\&.
                    \bigbreak \noindent 
                    \begin{cppcode}
                        void f(int&& x) { ... }

                        int x = 10;
                        f(10); // prvalue binds
                        f(std::move(x)) // xvalue binds
                    \end{cppcode}
                    \bigbreak \noindent 
                    \begin{cppcode}
                        void g(const int& x) { ... }
                        g(10); // prvalue binds
                        g(x); // lvalue binds
                    \end{cppcode}
                \item \textbf{Does std::move(x) move anything by itself? What does it do to the expression x?}: std::move(x) does not move data; it casts x to an rvalue reference, making the resulting expression an xvalue. A move may happen later if that expression is used to call a move constructor or move assignment operator.
                \item \textbf{Inside this function, is the expression s an lvalue or an rvalue? Why?}:
                    \bigbreak \noindent 
                    \begin{cppcode}
                        void f(std::string&& s) {
                            use(s);
                        }
                    \end{cppcode}
                    \bigbreak \noindent 
                    $s$ has rvalue-reference type, but the expression $s$ is an lvalue because it names an object. To pass it on as an rvalue, you would write std::move(s).
                    \bigbreak \noindent 
                    The rvalue-reference type matters mainly at the function boundary: it determines what arguments may be passed to f.
                    \bigbreak \noindent 
                    \begin{cppcode}
                        void f(std::string&& s);

                        std::string name = "Ada";

                        f(name);                 // Error: lvalue cannot bind to string&&
                        f(std::move(name));      // OK: xvalue
                        f(std::string{"Ada"});   // OK: prvalue
                    \end{cppcode}
                    \bigbreak \noindent 
                    Inside $f$, merely writing $s$ treats it as an lvalue. To take advantage of the fact that $f$ received an object that may be moved from, explicitly use std::move(s):
                \item \textbf{What are the deduced type $T$ and the resulting parameter type T\&\& in each call and explain how reference collapsing affects each result}: 
                    \bigbreak \noindent 
                    \begin{cppcode}
                        template<class T>
                        void f(T&& x);

                        std::string name = "Ada";

                        f(name);
                        f(std::string{"Ada"});
                    \end{cppcode}
                    \bigbreak \noindent 
                    For the lvalue
                    \bigbreak \noindent 
                    \begin{cppcode}
                    f(name);
                    \end{cppcode}
                    \bigbreak \noindent 
                    $T$ is deduced as std::string\&, referencing collapsing yields \& + \&\& = \&. Thus, the function argument is type std::string\&.
                    \bigbreak \noindent 
                    For the temporary:
                    \begin{cppcode}
                        f(std::string{"Ada"});
                    \end{cppcode}
                    \bigbreak \noindent 
                    $T$ is deduced as std::string. $T$\&\& therefore becomes std::string\&\&. This makes T\&\& a forwarding reference because $T$ is deduced from the argument.
                \item \textbf{Why should this wrapper use std::forward<T>(value) instead of std::move(value)? Assume consume has overloads for lvalues and rvalues}: 
                    \bigbreak \noindent 
                    \begin{cppcode}
                        template<class T>
                        void wrapper(T&& value) {
                            consume(/* value */);
                        }
                    \end{cppcode}
                    \begin{itemize}
                        \item If the original argument was an lvalue, T is U\&, so forwarding produces an lvalue.
                        \item If the original argument was an rvalue, T is U, so forwarding produces an xvalue.
                        \item std::move(value) would produce an xvalue in both cases and could incorrectly move from an original lvalue.
                    \end{itemize}
            \end{enumerate}
            \textbf{Note:} Consider
            \bigbreak \noindent 
            \begin{cppcode}
                template<typename T>
                void f(T&& arg) { ... }

                template<typename T>
                void g(T arg) { ... }

                int x{5};
                f(x); g(x);
            \end{cppcode}
            \bigbreak \noindent 
            In the call to $g$, the type is deduced as \texttt{int}, whereas in $f$, the type is deduced as \texttt{int}\&. The fundamental difference in behavior exists to prevent silent, unintentional copying of objects while enabling perfect forwarding.
            \bigbreak \noindent 
            In C++, template type deduction rules are explicitly designed around the signature of the function parameter. The compiler treats T arg and T\&\& arg with two entirely different sets of rules when it encounters an lvalue.
        \item \textbf{Implementation problem}: Implement forward\_to\_consume so it calls the appropriate overload:
            \bigbreak \noindent 
            \begin{cppcode}
                void consume(const std::string&) {
                    std::cout << "lvalue\n";
                }

                void consume(std::string&&) {
                    std::cout << "rvalue\n";
                }

                template<typename T>
                void forward_to_consume(T&& value) {
                    std::forward<T>(value);
                }

                int main() {
                    std::string name = "Ada";

                    forward_to_consume(name);                // Must print: lvalue
                    forward_to_consume(std::string{"Grace"}); // Must print: rvalue
                }
            \end{cppcode}
        \item \textbf{Dependent name qualifier}: A \textbf{dependent name} is a name that depends on a template parameter.
            \bigbreak \noindent 
            Because the compiler doesn't know what type will eventually replace the template parameter during compilation, it cannot look up what a dependent name actually refers to until the template is instantiated
            \bigbreak \noindent 
            \begin{cppcode}
                template <typename T>
                void process(T obj) {
                    int x = 10;        // 'int' and 'x' are Non-Dependent Names. They mean the same thing regardless of T.
                    T::value_type y;   // 'value_type' is a Dependent Name. What it means depends entirely on what T is.
                }
            \end{cppcode}
            \bigbreak \noindent 
            Dependent names generally fall into two categories, both of which require special keywords to help the compiler parse the code before instantiation.
            \bigbreak \noindent 
            \begin{enumerate}
                \item \textbf{Dependant types (requires typename)}
                    Because the compiler defaults to assuming nested names are member variables or functions, you must explicitly prefix it with typename.
                    \bigbreak \noindent 
                    \begin{cppcode}
                        template <typename T>
                        void print_size() {
                            // Compiler error without 'typename' because it assumes 'size_type' is a static variable
                            typename T::size_type size = 0; 
                        }
                    \end{cppcode}
                \item \textbf{Dependant templates (requires template)}: A member accessed through a dependent type might itself be a template
                    \bigbreak \noindent 
                    \begin{cppcode}
                        template<typename T>
                        void call(T& object) {
                            object.template convert<int>();
                        }
                    \end{cppcode}
                    \bigbreak \noindent 
                    The template keyword tells the compiler that convert is a member template and that <int> contains template arguments.
                    \bigbreak \noindent 
                    Without it:
                    \bigbreak \noindent 
                    \begin{cppcode}
                    object.convert<int>();
                    \end{cppcode}
                    \bigbreak \noindent 
                    the compiler might interpret < as the less-than operator because it does not yet know what members T has.
                    \bigbreak \noindent 
                    \begin{cppcode}
                        struct Converter {
                            template<typename U>
                            U convert() {
                                return U{};
                            }
                        };

                        template<typename T>
                        void use(T& object)
                        {
                            int result = object.template convert<int>();
                        }
                    \end{cppcode}
            \end{enumerate}
            \bigbreak \noindent 
            Use typename when a dependent qualified name represents a type. Use template when a member accessed through a dependent expression is a template.
        \item \textbf{Dependent base classes}: A base class can depend on a template parameter:
            \bigbreak \noindent 
            \begin{cppcode}
                template<typename T>
                struct Base {
                    void print() {}
                };

                template<typename T>
                struct Derived : Base<T> {
                    void run()
                    {
                        print();  // May fail
                    }
                };
            \end{cppcode}
            \bigbreak \noindent 
            Base<T> is a dependent base class. During the initial parsing of Derived, the compiler does not search dependent base classes for unqualified names.
            \bigbreak \noindent 
            You can qualify the call:
            \bigbreak \noindent 
            \begin{cppcode}
                template<typename T>
                struct Derived : Base<T> {
                    void run() {
                        this->print();
                    }
                };
            \end{cppcode}
            \bigbreak \noindent 
            Other valid approaches include:
            \bigbreak \noindent 
            \begin{cppcode}
            Base<T>::print();
            \end{cppcode}
            \bigbreak \noindent 
            or:
            \bigbreak \noindent 
            \begin{cppcode}
            using Base<T>::print;
            \end{cppcode}
            \bigbreak \noindent 
            The this-> expression is dependent because the type of this depends on T. Therefore, lookup for print can be delayed until instantiation.
        \item \textbf{Two-phase name lookup}: Templates are generally checked in two phases.
            \begin{enumerate}
                \item \textbf{Phase 1: Template definition}: The compiler parses the template and resolves names that do not depend on template parameters.
                    \bigbreak \noindent 
                    \begin{cppcode}
                        void helper();

                        template<typename T>
                        void function(T value)
                        {
                            helper();       // Non-dependent
                            value.work();   // Dependent
                        }
                    \end{cppcode}
                    \bigbreak \noindent 
                    helper() can be looked up immediately. The validity of value.work() depends on T.
                \item \textbf{Phase 2: Template instantiation}: When the template is used with a concrete argument, dependent names and expressions are resolved:
                    \bigbreak \noindent 
                    \begin{cppcode}
                        struct Worker {
                            void work() {}
                        };

                        function(Worker{});  // value.work() becomes valid
                        function(42);        // Error: int has no work()
                    \end{cppcode}
                    \bigbreak \noindent 
                    This delayed checking is one reason template errors often appear at the point where the template is instantiated.
            \end{enumerate}
        \item \textbf{Uses of T\&\&}: T\&\& has two closely related uses:
            \begin{enumerate}
                \item An ordinary rvalue reference binds to rvalues and enables move semantics, allowing resources such as a vector’s heap allocation to be transferred instead of copied.
                \item In a deduced template context, such as template<class T> void f(T\&\& x), it may be a forwarding reference. It can accept both lvalues and rvalues, and std::forward<T>(x) preserves the argument’s original value category.
            \end{enumerate}
        \item \textbf{ADL}: Argument-Dependent Lookup (ADL), also known as Koenig lookup, is a set of rules in C++ that allows the compiler to look up an unqualified function name within the namespaces of its arguments
            \bigbreak \noindent 
            Without ADL, you would be forced to explicitly qualify every single function or operator call with its full namespace
            \bigbreak \noindent 
            When you make a function call like foo(arg1, arg2), the compiler performs two lookup phases to form an overload set:
            \begin{enumerate}
                \item \textbf{Normal Unqualified Lookup:} The compiler searches the current local scope, enclosing scopes, and global namespaces
                \item \textbf{ADL Lookup:} The compiler examines the types of arg1 and arg2. It discovers their "associated namespaces" and "associated classes" and searches those namespaces as well
            \end{enumerate}
            ADL is skipped entirely if normal unqualified lookup finds any of the following
            \begin{enumerate}
                \item \textbf{A class member declaration:} If you call a function name that matches a member function inside a class method.
                \item \textbf{A block-scope function declaration:} A local function declaration inside a pair of braces \{ ... \} (excluding using declarations)
                \item \textbf{A non-function declaration:} If a local variable, function object (functor), or constant shares the same name as the function you are trying to call
            \end{enumerate}
            ADL is vital for writing generic template code, especially for customization points like std::swap. The standard idiom requires bringing std::swap into scope but calling it un-qualified:
            \bigbreak \noindent 
            \begin{cppcode}
                template <typename T>
                void doSomething(T& a, T& b) {
                    using std::swap; // Makes std::swap a fallback option
                    swap(a, b);      // Unqualified call triggers ADL!
                }
            \end{cppcode}
            \begin{itemize}
                \item If T is a user-defined type with a custom swap function defined in its own namespace, ADL will find and prioritize the custom version.
                \item If no custom swap exists, the compiler falls back to std::swap.
            \end{itemize}
            If you call swap(a, b) using an unqualified call without bringing std::swap into scope, ADL will fail for primitive types (like int, double, pointer) and types in the global namespace. This is because primitive types do not belong to any namespace, meaning ADL has no associated namespaces to search.
            \bigbreak \noindent 
            In C++20 and later, you no longer need to write this two-line ritual if you use Concepts and Ranges. The standard library introduced std::ranges::swap, which is a specialized function object (a "niebloid") that automatically handles ADL fallback logic under the hood:






    \end{itemize}


    \pagebreak 
    \unsect{Threads and Thread safety}
    \subsection{Threads and locks}
    \begin{itemize}
        \item \textbf{Creating a thread}
            \bigbreak \noindent 
            \begin{cppcode}
                #include <iostream>
                #include <thread>

                void worker(int id) {
                    std::cout << "Worker " << id << '\n';
                }

                int main() {
                    std::thread t(worker, 1);

                    t.join();
                }
            \end{cppcode}
            \bigbreak \noindent 
            The operating system schedules the new thread independently. You should not assume which thread runs first or how their instructions will be interleaved.
        \item \textbf{Join}: A std::thread is normally handled with join().
            \bigbreak \noindent 
            \begin{cppcode}
                std::thread t(do_work);
                t.join();
            \end{cppcode}
            \bigbreak \noindent 
            join() blocks the calling thread until t finishes. A thread object must not be destroyed while it is still joinable. Doing so calls std::terminate():
            \bigbreak \noindent 
            A thread must be joined to make the calling thread wait until the target thread finishes its work, prevent the main program from exiting too early, and safely clean up system resources
            \bigbreak \noindent 
            You can check its state:
            \bigbreak \noindent 
            \begin{cppcode}
            if (t.joinable())
                t.join();
            \end{cppcode}
        \item \textbf{Detached threads}:
            \bigbreak \noindent 
            \begin{cppcode}
            t.detach();
            \end{cppcode}
            \bigbreak \noindent 
            A detached thread continues independently. The std::thread object no longer controls it.
            \bigbreak \noindent 
            Detached threads are dangerous because the program must ensure that everything the thread accesses remains alive:
            \bigbreak \noindent 
            \begin{cppcode}
                void dangerous() {
                    int value = 42;

                    std::thread([&value] {
                        use(value);  // value might already be destroyed
                    }).detach();
                }
            \end{cppcode}
        \item \textbf{jthread}: C++20 introduced std::jthread. Its destructor automatically joins the thread.
            \bigbreak \noindent 
            \begin{cppcode}
                #include <thread>

                void worker() {
                    // Work
                }

                int main() {
                    std::jthread t(worker);
                } // Automatically joined
            \end{cppcode}
        \item \textbf{Passing arguments}: Arguments passed to std::thread are normally copied or moved into internal storage.
            \bigbreak \noindent 
            \begin{cppcode}
                void process(int value);

                int x = 10;
                std::thread t(process, x); // x is copied
            \end{cppcode}
            \bigbreak \noindent 
            To deliberately pass by reference, use std::ref:
            \bigbreak \noindent 
            \begin{cppcode}
                #include <functional>

                void increment(int& value) {
                    ++value;
                }

                int x = 10;
                std::thread t(increment, std::ref(x));
                t.join();
            \end{cppcode}
        \item \textbf{std::mutex}: A mutex provides mutual exclusion: only one thread can own it at a time.
            \bigbreak \noindent 
            \begin{cppcode}
                #include <mutex>

                int counter = 0;
                std::mutex counter_mutex;

                void increment()
                {
                    for (int i = 0; i < 100'000; ++i) {
                        counter_mutex.lock();
                        ++counter;
                        counter_mutex.unlock();
                    }
                }
            \end{cppcode}
            \bigbreak \noindent 
            The protected region is called a critical section. Do not normally call lock() and unlock() manually. If an exception or early return occurs before unlock(), the mutex remains locked
        \item \textbf{RAII locking}: C++ locking should normally use RAII: lock in a constructor, unlock in a destructor.
            \bigbreak \noindent 
            Use std::lock\_guard for a simple scope-based lock:
            \bigbreak \noindent 
            \begin{cppcode}
                void increment() {
                    for (int i = 0; i < 100'000; ++i) {
                        std::lock_guard<std::mutex> lock(counter_mutex);
                        ++counter;
                    }
                }
            \end{cppcode}
            \bigbreak \noindent 
            The mutex is automatically unlocked at the end of the scope.
            \bigbreak \noindent 
            You can shorten the critical section:
            \bigbreak \noindent 
            \begin{cppcode}
                void function() {
                    perform_unrelated_work();

                    {
                        std::lock_guard<std::mutex> lock(counter_mutex);
                        modify_shared_state();
                    }

                    perform_more_unrelated_work();
                }
            \end{cppcode}
            \bigbreak \noindent 
            Keeping critical sections short reduces contention.
            \bigbreak \noindent 
            When you initialize a std::lock\_guard, it calls .lock() on the underlying mutex (usually std::mutex). If another thread holds the lock, the operating system puts the requesting thread to sleep (context switch).
        \item \textbf{std::unique\_lock}: std::unique\_lock is more flexible:
            \bigbreak \noindent 
            \begin{cppcode}
                std::unique_lock<std::mutex> lock(counter_mutex);

                modify_shared_state();

                lock.unlock();
                perform_expensive_work();

                lock.lock();
                modify_shared_state_again();
            \end{cppcode}
            \bigbreak \noindent 
            It supports:
            \begin{itemize}
                \item Deferred locking
                \item Manual unlock/relock
                \item Ownership transfer
                \item Condition variables
                \item Timed mutex operations
            \end{itemize}
            It has slightly more state and overhead than std::lock\_guard.
        \item \textbf{std::scoped\_lock is particularly useful when acquiring several mutexes}: 
            \bigbreak \noindent 
            \begin{cppcode}
            std::scoped_lock lock(account_a.mutex, account_b.mutex);
            \end{cppcode}
            \bigbreak \noindent 
            It uses a deadlock-avoidance algorithm rather than naively locking them in the listed sequence.
        \item \textbf{try\_lock() and timed mutexes}: try\_lock() attempts to acquire a mutex without waiting:
            \bigbreak \noindent 
            \begin{cppcode}
                if (mutex.try_lock()) {
                    // Mutex acquired
                    use_shared_state();
                    mutex.unlock();
                } else {
                    // Do something else
                }
            \end{cppcode}
            \bigbreak \noindent 
            The primary difference is that std::mutex::lock() is a blocking operation, while std::mutex::try\_lock() is a non-blocking operation
            \bigbreak \noindent 
            RAII can still be used:
            \bigbreak \noindent 
            \begin{cppcode}
                std::unique_lock lock(mutex, std::try_to_lock);

                if (lock.owns_lock()) {
                    use_shared_state();
                }
            \end{cppcode}
        \item \textbf{Condition variables}: In C++, std::condition\_variable is a synchronization primitive that allows one thread to suspend execution (sleep) until another thread modifies a shared variable and signals the condition variable. This eliminates the need for inefficient busy-waiting loops
            \bigbreak \noindent 
            To use a condition variable, you need three pieces working together
            \begin{enumerate}
                \item \textbf{std::condition\_variable:} The signaling mechanism.
                \item \textbf{std::mutex:} Protects the shared state.
                \item \textbf{std::unique\_lock<std::mutex>:} Passed to the wait() function. A std::lock\_guard cannot be used for waiting because the condition variable must frequently unlock and relock the mutex while waiting
            \end{enumerate}
            \bigbreak \noindent 
            \begin{cppcode}
                #include <condition_variable>
                #include <mutex>
                #include <queue>

                std::queue<int> queue;
                std::mutex mutex;
                std::condition_variable available;
                bool finished = false;
            \end{cppcode}
            \bigbreak \noindent 
            \begin{cppcode}
            // Producer
            void produce(int value) {
                {
                    std::lock_guard lock(mutex);
                    queue.push(value);
                }

                available.notify_one();
            }
            \end{cppcode}
            \bigbreak \noindent 
            \begin{cppcode}
            // Consumer

            void consume() {
                for (;;) {
                    std::unique_lock lock(mutex);

                    available.wait(lock, [] {
                        return !queue.empty() || finished;
                    });

                    if (queue.empty() && finished) {
                        break;
                    }

                    int value = queue.front();
                    queue.pop();

                    lock.unlock();
                    process(value);
                }
            }
            \end{cppcode}
            \bigbreak \noindent 
            The predicate must be checked because:
            \begin{itemize}
                \item Wakeups can be spurious.
                \item Another consumer may take the item first.
                \item A notification means the state may have changed, not that the condition is definitely true.
            \end{itemize}
            \textbf{Note:} A spurious wakeup occurs when a thread wakes up from a waiting state (sleeping/blocked on a condition variable or monitor) without receiving any explicit signal, notification, or interrupt, and before the condition it is waiting for has actually been met.
            \bigbreak \noindent 
            The actual condition is the protected shared state, such as !queue.empty(). The condition variable is only the notification mechanism.
        \item \textbf{std::shared\_lock}: std::shared\_lock is a general-purpose, RAII-compliant shared lock ownership wrapper introduced in C++14. It is primarily designed for reader-writer scenarios, allowing multiple "reader" threads to concurrently access a shared resource while ensuring that "writer" threads get exclusive access
        \item \textbf{Reader-writer locking}: std::shared\_mutex permits:
            \begin{itemize}
                \item Multiple simultaneous readers
                \item One exclusive writer
            \end{itemize}
            std::shared\_mutex permits:
            \begin{itemize}
                \item Multiple simultaneous readers
                \item One exclusive writer
            \end{itemize}
            \begin{itemize}
                \item \textbf{Multiple Readers:} Multiple threads can hold a std::shared\_lock on the same mutex simultaneously
                \item \textbf{Exclusive Block:} If any thread holds an exclusive lock (e.g., via std::unique\_lock or std::lock\_guard), all std::shared\_lock requests will block until the exclusive lock is released
                \item \textbf{RAII Management:} It automatically calls lock\_shared() on construction and unlock\_shared() upon destruction (when it goes out of scope
            \end{itemize}
            \bigbreak \noindent 
            \begin{cppcode}
                #include <shared_mutex>
                #include <unordered_map>

                class Dictionary {
                    public:
                    int get(int key) const {
                        std::shared_lock lock(mutex_);
                        return values_.at(key);
                    }

                    void set(int key, int value) {
                        std::unique_lock lock(mutex_);
                        values_[key] = value;
                    }

                    private:
                        mutable std::shared_mutex mutex_;
                        std::unordered_map<int, int> values_;
                };
            \end{cppcode}
            \bigbreak \noindent 
            This helps primarily when reads are frequent, writes are uncommon, and the protected operations are substantial enough to justify the added complexity. An ordinary mutex may otherwise perform better.
        \item \textbf{One-time initialization}: Use std::call\_once when initialization must occur exactly once:
            \bigbreak \noindent 
            \begin{cppcode}
                #include <mutex>

                std::once_flag initialization_flag;

                void initialize()
                {
                    std::call_once(initialization_flag, [] {
                        initialize_library();
                    });
                }
            \end{cppcode}
            \bigbreak \noindent 
            Function-local static initialization is also thread-safe since C++11:
            \bigbreak \noindent 
            \begin{cppcode}
                Resource& resource() {
                    static Resource instance;
                    return instance;
                }
            \end{cppcode}
        \item \textbf{Semaphore}: A semaphore represents a number of available permits:
            \bigbreak \noindent 
            \begin{cppcode}
                #include <semaphore>

                std::counting_semaphore<4> slots(4);

                void worker() {
                    slots.acquire();
                    use_limited_resource();
                    slots.release();
                }
            \end{cppcode}
            \bigbreak \noindent 
            Unlike a mutex, a semaphore does not represent ownership of a critical section. It is useful for limiting concurrency or representing available resources.
            \bigbreak \noindent 
            A semaphore (whether a counting semaphore or a binary semaphore) does not represent ownership because it is a signaling mechanism, not a locking mechanism.
            \bigbreak \noindent 
            Unlike a mutex (mutual exclusion lock), which is explicitly designed to manage ownership of a critical section, a semaphore focuses strictly on managing a count of available resources.
            \begin{itemize}
                \item \textbf{The Rule:} In a mutex, only the exact thread that locked it can unlock it.
                \item \textbf{The Semaphore Difference:} Any thread can call signal() (or post/V) on a semaphore, even if it never called wait() (or pend/P) to acquire it. Because there is no restriction tying a specific thread to the resource, no concept of "ownership" exists.
            \end{itemize}
        \item \textbf{Latch}: A std::latch is a one-use countdown:
            \bigbreak \noindent 
            \begin{cppcode}
                #include <latch>
                #include <thread>
                #include <vector>

                std::latch completed(4);

                void worker() {
                    perform_work();
                    completed.count_down();
                }

                // Start four workers...
                completed.wait();
            \end{cppcode}
            \bigbreak \noindent 
            std::latch is a C++20 thread synchronization primitive that pauses one or more threads until a downward counter reaches zero
            \begin{itemize}
                \item \textbf{Single-Use:} Once the internal counter hits zero, the latch is done and cannot be reset or reused.
                \item \textbf{No Increment:} You cannot increase the counter after creation; you only count it down.
            \end{itemize}
            std::latch::wait blocks the calling thread until the internal counter of the std::latch reaches zero.
            \bigbreak \noindent 
            \begin{cppcode}
                #include <latches>
                #include <thread>
                #include <vector>

                // Initialize latch with count for 3 workers
                std::latch work_latch(3);

                for (int i = 0; i < 3; ++i) {
                    std::thread([&work_latch]() {
                        // Do work...
                        work_latch.count_down(); // Signal completion
                    }).detach();
                }

                work_latch.wait(); // Main thread blocks here until all 3 workers finish
            \end{cppcode}
        \item \textbf{std::barrier}: std::barrier is a C++20 thread-coordination mechanism that blocks a group of participating threads until a specific number of arrivals are reached, after which it can be reused for multiple phases. 
            \bigbreak \noindent 
             Unlike std::latch, which is single-use and destroys its counter value at zero, std::barrier automatically resets and can be used for repetitive coordination cycles
             \bigbreak \noindent 
             When all threads arrive at the synchronization point, the barrier executes a completion step (callable object) on one of the participating threads before unblocking the rest
             \bigbreak \noindent 
             \begin{cppcode}
                 #include <barrier>
                 int main() {
                     // 1. Define a completion function to run at the end of each phase
                     auto on_completion = []() noexcept {
                          ...
                     };

                     std::barrier sync_point(2, on_completion);

                     auto worker = [&](std::string name) {
                         ...

                         // Block until all 2 threads reach this point
                         sync_point.arrive_and_wait(); 

                         ...

                         // Reuse the same barrier for the next phase
                         sync_point.arrive_and_wait(); 

                         ...
                     };

                     std::vector<std::thread> threads;
                     threads.emplace_back(worker, "Thread A");
                     threads.emplace_back(worker, "Thread B");

                     for (auto& t : threads) {
                         t.join();
                     }

                     return 0;
                 }
             \end{cppcode}
            \item \textbf{std::async}: std::async is a function template in C++ (introduced in C++11) that runs a callable object-like a function, lambda, or functor-asynchronously or lazily, returning a std::future that will eventually hold the result. It abstracts away manual thread management, automatically handling execution and allowing you to retrieve return values or handle exceptions across thread boundaries safely
            \item \textbf{std::promise}: A std::promise represents a commitment made by a worker thread to provide a result at some point.  The thread fulfills this contract by calling .set\_value().
            \item \textbf{std::future}: A std::future represents a placeholder for a value that does not exist yet, but will exist in the future.
                \bigbreak \noindent 
                When you call .get(), you are looking into the "future." If the future hasn't arrived yet, you sit and wait for it.
            \item \textbf{std::promise with std::future}: In C++, std::promise and std::future are a pair of class templates introduced in C++11 to pass values, exceptions, or signals between threads without explicitly managing mutexes or condition variables. They establish a one-way, single-use communication channel where one thread acts as the producer and another acts as the consumer
                \bigbreak \noindent 
                \begin{cppcode}
                    void calculate(std::promise<int> result)
                    {
                        try {
                            result.set_value(42);
                        } catch (...) {
                            result.set_exception(std::current_exception());
                        }
                    }

                    int main()
                    {
                        std::promise<int> promise;
                        std::future<int> future = promise.get_future();

                        std::thread t(calculate, std::move(promise));

                        int result = future.get();
                        t.join();
                    }
                \end{cppcode}
                \bigbreak \noindent 
                Promises are useful when the producer and consumer need a direct one-time communication channel.



    \end{itemize}


    \pagebreak 
    \subsection{Atomics}
    \begin{itemize}
        \item \textbf{std::atomic<T>}: std::atomic<T> provides thread-safe access to a shared value without requiring an explicit mutex for each operation.        
            \bigbreak \noindent 
            Its main purpose is preventing a data race when multiple threads read or modify the same variable.
            \bigbreak \noindent 
            Without an atomic variable, this program has undefined behavior
            \bigbreak \noindent 
            \begin{cppcode}
                #include <thread>

                int counter = 0;

                void increment()
                {
                    for (int i = 0; i < 100'000; ++i) {
                        ++counter;
                    }
                }
            \end{cppcode}
            \bigbreak \noindent 
            ++counter is a read-modify-write operation. Two threads can interfere with each other. Instead, we can use std::atomic
            \bigbreak \noindent 
            \begin{cppcode}
                #include <atomic>
                #include <iostream>
                #include <thread>

                std::atomic<int> counter{0};

                void increment()
                {
                    for (int i = 0; i < 100'000; ++i) {
                        ++counter;
                    }
                }

                int main()
                {
                    std::thread first(increment);
                    std::thread second(increment);

                    first.join();
                    second.join();

                    std::cout << counter.load() << '\n';  // 200000
                }
            \end{cppcode}
            \bigbreak \noindent 
            Every increment is now indivisible from the perspective of other threads.
        \item \textbf{Operations}: 
            \bigbreak \noindent 
            \begin{cppcode}
                std::atomic<int> value{10};

                value.store(20);          // Atomic write
                int x = value.load();     // Atomic read

                value.fetch_add(5);       // Add 5; returns the old value
                value.fetch_sub(2);       // Subtract 2; returns the old value

                value++;                  // Atomic increment
                value--;                  // Atomic decrement
            \end{cppcode}
        \item \textbf{Cannot be copied}: Atomic objects generally cannot be copied. You copy the contained value explicitly
        \item \textbf{Atomic flags}: A common use is communicating a simple condition between threads:
            \bigbreak \noindent 
            \begin{ccode}
                #include <atomic>
                #include <thread>

                std::atomic<bool> stop_requested{false};

                void worker()
                {
                    while (!stop_requested.load()) {
                        // Perform work
                    }
                }

                int main()
                {
                    std::thread thread(worker);

                    // Later:
                    stop_requested.store(true);

                    thread.join();
                }
            \end{ccode}
            \bigbreak \noindent 
            Atomicity ensures that accessing stop\_requested does not create a data race. If stop\_requested were a normal bool, the worker thread would read it while the main thread writes it. One thread writes the object while another thread reads it, producing a data race. In C++, a data race results in undefined behavior.
            \bigbreak \noindent 
            Because ordinary C++ code has no defined data races, the compiler may assume another thread does not modify stop\_requested. It could effectively transform this:
            \bigbreak \noindent 
            \begin{cppcode}
                while (!stop_requested) {
                }
            \end{cppcode}
            \bigbreak \noindent 
            into something resembling:
            \bigbreak \noindent 
            \begin{cppcode}
                bool cached = stop_requested;

                while (!cached) {
                }
            \end{cppcode}
            \bigbreak \noindent 
            The worker might therefore loop forever even after the main thread assigns true.
            \bigbreak \noindent 
            Each CPU core has registers, caches, and store buffers. A normal write does not provide the synchronization required by the C++ memory model to make another thread observe the change correctly.
            \bigbreak \noindent 
            Modern hardware has cache-coherence mechanisms, but those mechanisms alone do not make racy C++ code valid. The compiler and processor are both allowed to reorder ordinary operations when no synchronization is present.
            \bigbreak \noindent 
            On many systems, reading or writing a single byte-sized bool happens as one hardware operation. That only means the reader probably will not observe a partially written value.
            \bigbreak \noindent 
            Thread safety requires more than preventing a torn read. It also requires:
            \bigbreak \noindent 
            \begin{itemize}
                \item Defined concurrent access
                \item Appropriate visibility between threads
                \item Restrictions on compiler and processor reordering
            \end{itemize}
            A normal bool provides none of these guarantees.
        \item \textbf{Torn read}: A torn read is a concurrency bug that occurs when a thread reads a value from memory while another thread is simultaneously writing to it, resulting in the reader seeing a partially updated, corrupted value
        \item \textbf{Cache coherence}: Cache coherence is the practice of making sure that all copies of shared data stored in different local caches remain identical in a multi-core or multi-processor system
        \item \textbf{Compare-and-exchange}: Compare-and-exchange is central to lock-free algorithms. It means if the atomic currently contains the expected value, replace it with the desired value. Otherwise, update the expected variable with the actual value.
            \bigbreak \noindent 
            \begin{cppcode}
                std::atomic<int> value{10};

                int expected = 10;

                bool changed = value.compare_exchange_strong(expected, 20);
            \end{cppcode}
        \item \textbf{strong vs weak compare exchange}: compare\_exchange\_weak and compare\_exchange\_strong perform a conditional atomic update called compare-and-exchange, often abbreviated CAS.
            \bigbreak \noindent 
            Conceptually, CAS performs this indivisible operation:
            \bigbreak \noindent 
            \begin{cppcode}
                if (atomic_value == expected) {
                    atomic_value = desired;
                    return true;
                } else {
                    expected = atomic_value;
                    return false;
                }
            \end{cppcode}
            \bigbreak \noindent 
            No other thread can interrupt the comparison and replacement as two separate steps.
            \bigbreak \noindent 
            Consider
            \bigbreak \noindent 
            \begin{cppcode}
                int current = value.load();
                value.store(current + 1);
            \end{cppcode}
            \bigbreak \noindent 
            Two threads could perform this sequence:
            \begin{center}
                \begin{tabularx}{\textwidth}{@{}XXXX@{}}
                    \toprule
                    \textbf{Step}&	\textbf{Thread A}&	\textbf{Thread B}& \textbf{Atomic value} \\
                    \midrule
                    1	&Loads 10		&10\\[2ex]
                    2	&	Loads 10	&10\\[2ex]
                    3	&Stores 11		&11\\[2ex]
                    4	&	Stores 11	&11 \\
                    \bottomrule
                \end{tabularx}
            \end{center}
            One increment is lost. CAS allows each thread to store the new value only if the old value has not changed since it was loaded
            \bigbreak \noindent 
            The fundamental difference is that compare\_exchange\_weak may report failure even when the atomic object equals expected. This is called a \textbf{spurious failure}
            \bigbreak \noindent 
            A \textbf{spurious failure} in a compare-and-exchange (or Compare-and-Swap / CAS) operation occurs when the function returns false (indicating failure) even though the current value in memory exactly matches the expected value. Instead of executing the swap, the operation acts as if the values did not match
            \bigbreak \noindent 
            On many RISC architectures (like ARM, PowerPC, or MIPS), atomic operations are implemented using a pair of instructions: Load-Link (LL) and Store-Conditional (SC)
            \bigbreak \noindent 
            The LL instruction tracks the memory address. The SC instruction only succeeds if no changes or specific hardware events occurred at that address since the LL instruction.
            \bigbreak \noindent 
            Context switches, system interrupts, or even basic CPU cache eviction events can cause the CPU to "lose the reservation" on that memory address. As a result, the SC instruction fails naturally, inducing a spurious failure even if no other thread changed the underlying value
            \bigbreak \noindent 
            Strong fails only when the comparison fails, it is suitable for a one-time attempt.
            \bigbreak \noindent 
            Weak may fail because
            \begin{itemize}
                \item The atomic value did not equal expected; or
                \item The atomic value did equal expected, but the implementation produced a spurious failure.
            \end{itemize}
            Using compare\_exchange\_strong in a retry loop is usually considered a performance anti-pattern because compare\_exchange\_weak is more efficient on many CPU architectures.
            \bigbreak \noindent 
            While compare\_exchange\_strong guarantees it will only fail if the value actually changed, this guarantee comes with a hardware cost that is unnecessary inside a loop.
            \bigbreak \noindent 
            So, weak is normally used in a loop.
        \item \textbf{Memory ordering}: Atomic operations have two related responsibilities:
            \begin{enumerate}
                \item Making an individual operation indivisible.
                \item Controlling how memory operations become visible between threads.
            \end{enumerate}
            Memory ordering defines how compilers and CPU hardware are allowed to reorder read and write operations across multiple threads
            \begin{itemize}
                \item \textbf{Single-threaded code:} Compilers and CPUs change instruction order to make code run faster. This reordering is invisible because the final result stays the same. 
                \item \textbf{Multi-threaded code:} Other threads can spot these reordered operations. Without rules, one thread might see updates in a completely different order than another thread wrote them, breaking synchronization
            \end{itemize}
        \item \textbf{Store-Buffer Forwarding (Store-Load Reordering)}: Imagine two threads running on two different CPU cores, sharing two variables (X and Y), both starting at 0.
            \begin{center}
                \begin{tabularx}{\textwidth}{@{}XX@{}}
                    \toprule
                    \textbf{Thread 1 (Core 1)}& \textbf{Thread 2 (Core 2)} \\
                    \midrule
                    1. Write 1 to X &3. Write 1 to Y  \\[2ex]
                    2. Read Y &4. Read X \\
                    \bottomrule
                \end{tabularx}
            \end{center}
            Logically, it should be impossible for both threads to read 0. At least one thread has to write its 1 before the other thread reads.
            \bigbreak \noindent 
            Yet, without memory barriers, both reads can return 0.
            \bigbreak \noindent 
            Writing data to the main RAM (or even the L1 cache) is slow. To speed things up, each CPU core has a tiny, ultra-fast temporary holding zone called a Store Buffer.
            \begin{enumerate}
                \item Core 1 executes Write 1 to X. Instead of waiting for X to travel to cache, Core 1 drops the write into its local Store Buffer and immediately moves on.
                \item Core 1 executes Read Y. It checks its own Store Buffer (nothing for Y there) and reads Y from the cache. Because Core 2 hasn't flushed its own buffer yet, Y is still 0.
                \item Core 2 does the exact same thing simultaneously. It drops Write 1 to Y into its own Store Buffer and immediately executes Read X, which returns 0.
                \item Finally, both cores flush their Store Buffers out to the shared cache.
            \end{enumerate}
            To stop the CPU from doing this, you must use a Memory Barrier (like C++'s std::memory\_order\_seq\_cst). This instruction forces the CPU core to completely empty its local Store Buffer and commit the data to shared cache before it is allowed to perform the subsequent Read.
        \item \textbf{Memory barrier}:  A memory barrier (also called a memory fence) is a specific CPU instruction that forces the processor and the compiler to maintain a strict order of memory operations. It acts as a wall that prevents reads and writes from crossing past it.
            \begin{center}
                \begin{tabularx}{\textwidth}{@{}XXX@{}}
                    \toprule
                    \textbf{Barrier Type}& \textbf{What It Does}& \textbf{Common Use Case} \\
                    \midrule
                    LoadLoad & Ensures older reads finish before newer reads can start. & Preventing a thread from reading stale data. \\[2ex]
                    StoreStore &Ensures older writes are visible to cache before newer writes happen. &Guaranteeing a payload is written before turning on a "Ready" flag. \\[2ex]
                    LoadStore & Ensures older reads finish before newer writes can happen. &Preventing early overwrites of data. \\[2ex]
                    StoreLoad &Ensures older writes are completely flushed before newer reads start. & The heaviest barrier. Required to prevent the store-buffer shortcut described earlier. \\
                    \bottomrule
                \end{tabularx}
            \end{center}
        \item \textbf{More on memory ordering}: Memory ordering controls how operations surrounding an atomic operation may be observed by other threads.
            \bigbreak \noindent 
            Every std::atomic operation is atomic regardless of its memory order. The ordering argument determines whether that operation also establishes ordering and visibility for other memory accesses.
            \bigbreak \noindent 
            Consider
            \bigbreak \noindent 
            \begin{cppcode}
                int data = 0;
                std::atomic<bool> ready{false};

                void producer() {
                    data = 42;
                    ready.store(true);
                }

                void consumer() {
                    while (!ready.load()) { }

                    std::cout << data << '\n';
                }
            \end{cppcode}
            \bigbreak \noindent 
            The intent is:
            \begin{itemize}
                \item Producer writes data.
                \item Producer sets ready.
                \item Consumer observes ready.
                \item Consumer reads data.
            \end{itemize}
            For this to work, the atomic flag must do more than safely transfer a Boolean value. It must also ensure that the write to data becomes visible before the consumer reads it. Memory ordering defines that relationship.
        \item \textbf{Program order is not necessarily observation order}: Within one thread, the source code says:
            \bigbreak \noindent 
            \begin{cppcode}
                data = 42;
                ready = true;
            \end{cppcode}
            \bigbreak \noindent 
            But several layers may reorder or delay operations:
            \begin{itemize}
                \item The compiler may rearrange instructions.
                \item The CPU may execute or retire memory operations out of order.
                \item Stores may temporarily remain in a store buffer.
                \item Different cores may observe stores at different times.
            \end{itemize}
            These transformations are allowed when they do not change the observable behavior of a correctly synchronized C++ program.
            \bigbreak \noindent 
            Atomic memory orders constrain these transformations.
            \bigbreak \noindent 
            Importantly, the C++ memory model is the authority. Hardware cache coherence alone does not make unsynchronized C++ code valid.
            \begin{itemize}
                \item \textbf{Sequenced-before:} ordering inside one thread.
                \item \textbf{Synchronizes-with:} ordering established between threads through synchronization operations.
                \item \textbf{Happens-before:} the combined ordering that makes earlier effects visible and prevents data races.
            \end{itemize}
        \item \textbf{Global order}: A global order in thread memory models is a single, total sequence of operations that every thread in a program agrees on and observes in the exact same way
            \bigbreak \noindent 
             If Thread A performs an operation before Thread B in this global sequence, every other thread in the system will witness the events in that exact same chronological order
             \bigbreak \noindent 
             CPUs and compilers cannot reorder operations across different threads in a way that violates this shared global sequence
             \bigbreak \noindent 
             It mimics a simple single-threaded program where instructions happen one after another, making concurrent code easy to reason about
        \item \textbf{Even more on memory ordering}: std::memory\_order controls two separate properties of an atomic operation:
            \begin{enumerate}
                \item The atomic operation itself: whether other threads can observe it partially.
                \item The ordering of other memory operations around it: whether ordinary reads and writes become visible across threads.
            \end{enumerate}
            Every std::atomic operation is atomic regardless of its memory order. The difference is synchronization.
            \begin{center}
                \begin{tabularx}{\textwidth}{@{}XXXXX@{}}
                    \toprule
                    \textbf{Order} & \textbf{Atomicity} & \textbf{Acquires earlier writes} & \textbf{Publishes earlier writes} & \textbf{Global order} \\ 
                    \midrule
                    `relaxed` & Yes & No & No & No \\ [2ex]
                    `consume` & Yes & Dependency-based historically & No & No \\ [2ex]
                    `acquire` & Yes & Yes & No & No \\ [2ex]
                    `release` & Yes & No & Yes & No \\ [2ex]
                    `acq\_rel` & Yes & Yes & Yes & No \\[2ex]
                    `seq\_cst` & Yes & Yes & Yes & Yes, among sequentially consistent operations \\ 
                    \bottomrule
                \end{tabularx}
            \end{center}
            Note that modern spelling is 
            \bigbreak \noindent 
            \begin{cppcode}
                std::memory_order::relaxed
                std::memory_order::acquire
                std::memory_order::release
                std::memory_order::acq_rel
                std::memory_order::seq_cst
            \end{cppcode}
            \bigbreak \noindent 
            The older spellings, such as std::memory\_order\_relaxed, are also commonly seen.
        \item \textbf{Release and acquire}: 
            \bigbreak \noindent 
            \begin{cppcode}
                int data = 0;
                std::atomic<bool> ready{false};
            \end{cppcode}
            \bigbreak \noindent 
            Thread 1: 
            \bigbreak \noindent 
            \begin{cppcode}
                data = 42;
                ready.store(true, std::memory_order::release);
            \end{cppcode}
            \bigbreak \noindent 
            Thread 2:
            \bigbreak \noindent 
            \begin{cppcode}
                while (!ready.load(std::memory_order::acquire)) {
                }

                std::cout << data;  // Guaranteed to print 42
            \end{cppcode}
            \begin{enumerate}
                \item data = 42 sequenced-before
                \item release store of true synchronizes-with
                \item acquire load that reads true sequenced-before
                \item reading data
            \end{enumerate}
            Therefore, writing data happens-before reading data.
            \bigbreak \noindent 
            The critical condition is that the acquire load must observe the value from the release store. Merely using release and acquire somewhere in the program is insufficient.
            \bigbreak \noindent 
            if the load reads false, the consumer keeps waiting. When it eventually reads the producer’s true, the release/acquire synchronization is established.
            \bigbreak \noindent 
            “Sequenced-before” describes ordering within one thread. Because statement A appears before statement B in the producer thread, A is sequenced-before B:
            \bigbreak \noindent 
            The release store effectively publishes everything that the producer did before it. Therefore, the write to data is part of what is being published.
        \item \textbf{memory\_order\_relaxed}: Relaxed ordering guarantees atomicity and a per-object modification order, but it creates no synchronization for other memory.
            \bigbreak \noindent 
            This is appropriate when you only need an accurate counter. Concurrent increments cannot overwrite each other.
            \bigbreak \noindent 
            However, relaxed ordering cannot normally publish surrounding data:
            \bigbreak \noindent 
            \begin{cppcode}
                int data = 0;
                std::atomic<bool> ready{false};

                // Thread 1
                data = 42;
                ready.store(true, std::memory_order::relaxed);

                // Thread 2
                if (ready.load(std::memory_order::relaxed)) {
                    std::cout << data;  // Data race: undefined behavior
                }
            \end{cppcode}
            \bigbreak \noindent 
            The atomic ready variable is safe, but nothing connects the ordinary write to data with the ordinary read from data.
        \item \textbf{memory\_order\_release}: Release ordering is used on stores and read-modify-write operations.
            \bigbreak \noindent 
            \begin{cppcode}
            ready.store(true, std::memory_order::release);
            \end{cppcode}
            \bigbreak \noindent 
            Its purpose is publication: operations sequenced before the release can become visible to a thread that performs a matching acquire operation.
            \bigbreak \noindent 
            Valid operations are 
            \bigbreak \noindent 
            \begin{cppcode}
                atomic.store(value, std::memory_order::release);
                atomic.exchange(value, std::memory_order::release);
                atomic.fetch_add(value, std::memory_order::release);
            \end{cppcode}
            \bigbreak \noindent 
            This is invalid:
            \bigbreak \noindent 
            \begin{cppcode}
            atomic.load(std::memory_order::release);  // Invalid
            \end{cppcode}
            \bigbreak \noindent 
            A load does not publish a value, so release ordering does not make sense for a pure load.
        \item \textbf{memory\_order\_acquire}: Acquire ordering is used on loads and read-modify-write operations.
            \bigbreak \noindent 
            \begin{cppcode}
            bool value = ready.load(std::memory_order::acquire);
            \end{cppcode}
            \bigbreak \noindent 
            If this operation reads from a corresponding release operation, everything published before that release becomes visible after the acquire.
            \bigbreak \noindent 
            An acquire load does not publish preceding operations to other threads. It receives information published by another thread.
            \bigbreak \noindent 
            A pure store does not read a value, so it cannot acquire anything.
        \item \textbf{memory\_order\_acq\_rel}: Acquire-release combines both directions:
            \begin{itemize}
                \item Acquire the writes published by another thread.
                \item Release the current thread’s preceding writes.
            \end{itemize}
            It is intended for atomic read-modify-write operations:
            \bigbreak \noindent 
            \begin{cppcode}
                atomic.fetch_add(1, std::memory_order::acq_rel);
                atomic.exchange(value, std::memory_order::acq_rel);
                atomic.compare_exchange_weak(
                    expected,
                    desired,
                    std::memory_order::acq_rel
                );
            \end{cppcode}
            \bigbreak \noindent 
            A read-modify-write operation both reads the old value and publishes a new value, so it can meaningfully perform both functions.
            \bigbreak \noindent 
            \begin{cppcode}
                std::atomic<int> state{0};

                int old = state.fetch_add(1, std::memory_order::acq_rel);
            \end{cppcode}
            \bigbreak \noindent 
            The operation can:
            \begin{itemize}
                \item Acquire information associated with the value it reads.
                \item Release preceding work through the new value it writes.
            \end{itemize}
            \bigbreak \noindent 
            This is invalid for pure loads and stores
            \bigbreak \noindent 
            \begin{cppcode}
                atomic.load(std::memory_order::acq_rel);   // Invalid
                atomic.store(1, std::memory_order::acq_rel); // Invalid
            \end{cppcode}
            \bigbreak \noindent 
            Use acquire for the load and release for the store.
        \item \textbf{memory\_order\_seq\_cst}: Sequential consistency provides acquire/release behavior plus a single total order over all sequentially consistent atomic operations.
            \bigbreak \noindent 
            \begin{cppcode}
                atomic.store(1);  // seq_cst
                int x = atomic.load();  // seq_cst
            \end{cppcode}
            The program behaves as though all seq\_cst operations were placed into one global sequence that every thread agrees upon. The standard requires a single total order for these operations
            \bigbreak \noindent 
            Consider
            \bigbreak \noindent 
            \begin{cppcode}
                std::atomic<int> x{0};
                std::atomic<int> y{0};

                int r1;
                int r2;
            \end{cppcode}
            \bigbreak \noindent 
            \begin{cppcode}
            // Thread 1
            x.store(1, std::memory_order::seq_cst);
            r1 = y.load(std::memory_order::seq_cst);

            // Thread 2
            y.store(1, std::memory_order::seq_cst);
            r2 = x.load(std::memory_order::seq_cst);
            \end{cppcode}
            \bigbreak \noindent 
            With exclusively sequentially consistent operations, this result is forbidden:
            \bigbreak \noindent 
            \begin{cppcode}
            r1 == 0 && r2 == 0
            \end{cppcode}
            \bigbreak \noindent 
            To obtain both zeros, each load would need to appear before the other thread’s store, producing an impossible cycle in the global sequentially consistent order.
        \item \textbf{Release/acquire synchronization is conditional}: An acquire load synchronizes with a release store only when the load reads the value written by that store.
            \bigbreak \noindent 
            So, with
            \bigbreak \noindent 
            \begin{cppcode}
                x.store(1, std::memory_order::release);
                r1 = y.load(std::memory_order::acquire);

                // Thread 2
                y.store(1, std::memory_order::release);
                r2 = x.load(std::memory_order::acquire);

                // r1 == r2 == 0 possible
            \end{cppcode}
            \bigbreak \noindent 
            For thread 1s load
            \bigbreak \noindent 
            \begin{cppcode}
                r1 = y.load(std::memory_order::acquire);
            \end{cppcode}
            \begin{itemize}
                \item If it reads 1, it read from Thread 2’s release store, so synchronization occurs.
                \item If it reads the original 0, it did not read from Thread 2’s release store, so no synchronization occurs.
            \end{itemize}
            The same applies to Thread 2’s load of x. Therefore, if both loads read the initial values, neither acquire load read from the other thread’s release store
            \bigbreak \noindent 
            \textbf{Note:} The standard explicitly restricts stores to relaxed, release, or sequentially consistent ordering, and loads to relaxed, acquire, or sequentially consistent ordering
        \item \textbf{Compare-exchange success and failure orders}: Compare-exchange often takes two orders:
            \bigbreak \noindent 
            \begin{cppcode}
                value.compare_exchange_weak(
                    expected,
                    desired,
                    std::memory_order_acq_rel,  // Success
                    std::memory_order_acquire   // Failure
                );
            \end{cppcode}
            \bigbreak \noindent 
            On success, it reads the old value and writes the new one, so acq\_rel is possible.
            \bigbreak \noindent 
            On failure, it only reads the actual value into expected. Because no write occurred, the failure order cannot contain release semantics. Therefore, release and acq\_rel are invalid failure orders.
            \bigbreak \noindent 
            A common pattern is 
            \bigbreak \noindent 
            \begin{cppcode}
                while (!head.compare_exchange_weak(
                    expected,
                    desired,
                    std::memory_order_release,
                    std::memory_order_relaxed
                )) {
                    // expected was updated
                }
            \end{cppcode}
            \bigbreak \noindent 
            The successful operation publishes data. Failure merely retrieves the current value and does not need synchronization.
        \item \textbf{Atomics do not make an entire algorithm atomic}: Suppose two atomic variables must remain consistent with each other:
            \bigbreak \noindent 
            \begin{cppcode}
                std::atomic<int> account_a{100};
                std::atomic<int> account_b{100};

                void transfer() {
                    account_a.fetch_sub(50);
                    account_b.fetch_add(50);
                }
            \end{cppcode}
            \bigbreak \noindent 
            Each operation is atomic, but the complete transfer is not. Another thread could observe:
            \bigbreak \noindent 
            \begin{cppcode}
                account_a == 50
                account_b == 100
            \end{cppcode}
            \bigbreak \noindent 
            If both values must change as one transaction, use a mutex:
            \bigbreak \noindent 
            \begin{cppcode}
                #include <mutex>

                int account_a = 100;
                int account_b = 100;
                std::mutex accounts_mutex;

                void transfer() {
                    std::lock_guard lock(accounts_mutex);

                    account_a -= 50;
                    account_b += 50;
                }
            \end{cppcode}
            \bigbreak \noindent 
            Atomics protect individual atomic operations. They do not automatically protect larger invariants involving several operations or objects.
        \item \textbf{Lock-free does not always mean faster}: An atomic type may be implemented directly using hardware instructions, or the implementation may internally use locks.
            \bigbreak \noindent 
            You can inspect this for a type
            \bigbreak \noindent 
            \begin{cppcode}
                std::atomic<long long> value{0};

                if (value.is_lock_free()) {
                    // This atomic object is implemented without locks.
                }
            \end{cppcode}
            \bigbreak \noindent 
            For a compile-time check when available:
            \bigbreak \noindent 
            \begin{cppcode}
            static_assert(std::atomic<int>::is_always_lock_free);
            \end{cppcode}
            \bigbreak \noindent 
            Simple types such as bool, pointers, and properly aligned machine-sized integers are commonly lock-free, but portable code should not assume this unless the implementation guarantees it.
        \item \textbf{Atomic pointers}: Pointers can also be atomic:
            \bigbreak \noindent 
            \begin{cppcode}
                int values[10];

                std::atomic<int*> pointer{values};

                pointer.fetch_add(1);  // Advances by one int
                int* current = pointer.load();
            \end{cppcode}
            \bigbreak \noindent 
            Atomic pointers are important in lock-free data structures, although safe memory reclamation makes those structures substantially more complicated than they initially appear.
        \item \textbf{C++20 waiting}: C++20 allows a thread to wait for an atomic value to change without manually spinning continuously:
            \bigbreak \noindent 
            \begin{cppcode}
            std::atomic<bool> ready{false};

            void consumer() {
                ready.wait(false);

                // ready is no longer false
            }

            void producer() {
                ready.store(true);
                ready.notify_one();
            }
            \end{cppcode}



    \end{itemize}

    \pagebreak 
    \unsect{The Embedded Memory Model}
    \begin{itemize}
        \item \textbf{Intro}: Embedded C++ uses the same language-level memory model as desktop C++, but embedded systems usually adopt a different allocation policy
            \bigbreak \noindent 
            Allocate everything at startup or compile time, then avoid dynamic allocation during normal operation.
            \bigbreak \noindent 
            Heap use is discouraged in embedded software because dynamic memory allocation introduces unpredictable behavior, memory fragmentation, and fatal crash risks that violate safety-critical requirement
            \bigbreak \noindent 
            Functions like malloc() or new take variable amounts of time to search for free memory blocks, which ruins real-time deadlines where timing must be exact
            \bigbreak \noindent 
            Allocating and freeing differently sized blocks over time breaks RAM into small, unusable gaps. A system can report plenty of free memory overall yet still fail a new allocation request
            \bigbreak \noindent 
            Small embedded devices have strict RAM limits. Memory leaks or unexpected allocation spikes cause the heap to run out, leading to hard faults or system resets
            \bigbreak \noindent 
            The stack and the heap grow toward each other in RAM; when they collide, the system crashes
        \item \textbf{Typical MCU memory layout}: A bare-metal firmware image commonly uses memory like this:
            \begin{verbatim}
                Flash
                |-- Vector table
                |-- Machine code (.text)
                |-- Constants (.rodata)
                |-- Initial values for global variables

                SRAM
                |-- Initialized globals/statics (.data)
                |-- Zero-initialized globals/statics (.bss)
                |-- Heap         usually grows upward
                |-- Unused RAM
                |-- Stack        usually grows downward
            \end{verbatim}
            \bigbreak \noindent 
            Code and constant data normally reside in flash, while mutable data, stacks, and any heap reside in SRAM.
        \item \textbf{Static-storage objects}: Objects can be allocated globally or as static objects. Their storage is reserved before the program begins normal execution. There is no runtime allocation and no fragmentation.
            \bigbreak \noindent 
            The downside is that these objects exist for the firmware’s entire lifetime and may complicate initialization or testing if overused.
        \item \textbf{Automatic stack allocation}: Short-lived objects are normally local variables:
            \bigbreak \noindent 
            \begin{cppcode}
                void process_message() {
                    std::array<std::byte, 64> message{};
                    Parser parser{message};

                    parser.process();
                }
            \end{cppcode}
            \bigbreak \noindent 
            The compiler adjusts the stack pointer when the function begins and restores it when the function returns. There is no allocator involved, so allocation is fast and deterministic. 
            \bigbreak \noindent 
            However, MCU stacks are limited. Avoid putting very large buffers on the stack:
            \bigbreak \noindent 
            Large buffers are generally better placed in static storage or a deliberately managed memory region.
        \item \textbf{Fixed-capacity containers}: Instead of containers that grow dynamically, embedded code often uses containers whose maximum capacity is known at compile time:
            \bigbreak \noindent 
            \begin{cppcode}
                template<typename T, std::size_t Capacity>
                class StaticVector {
                    std::array<T, Capacity> storage_;
                    std::size_t size_ = 0;

                    public:
                    bool push_back(const T& value)
                    {
                        if (size_ == Capacity)
                        return false;

                        storage_[size_++] = value;
                        return true;
                    }
                };
            \end{cppcode}
            This provides vector-like behavior without heap allocation.
            \bigbreak \noindent 
        \item \textbf{Fixed-size memory pools}: When objects must be created and destroyed dynamically, a fixed pool is often preferable to a general-purpose heap:
            \bigbreak \noindent 
            \begin{cppcode}
                struct Message {
                    std::uint32_t id;
                    std::array<std::byte, 32> payload;
                };

                std::array<Message, 16> message_pool;
            \end{cppcode}
            \bigbreak \noindent 
            A pool maintains a collection of equally sized slots. Allocation selects a free slot, and deallocation returns it.
            \begin{itemize}
                \item Fixed maximum memory usage
                \item No external fragmentation
                \item Usually constant or bounded allocation time
                \item Straightforward failure behavior
                \item Easier testing of exhaustion
            \end{itemize}
        \item \textbf{Arena allocation}: An arena owns a fixed region and dispenses memory sequentially:
            \bigbreak \noindent 
            \begin{cppcode}
                class Arena {
                    std::byte* current_;
                    std::byte* end_;

                    public:
                    Arena(void* memory, std::size_t size)
                        : current_{static_cast<std::byte*>(memory)},
                        end_{current_ + size} {}

                    void* allocate(std::size_t bytes);
                };
            \end{cppcode}
            \bigbreak \noindent 
            Individual objects usually are not freed. The complete arena is reset when all objects allocated from it are no longer needed.
        \item \textbf{New on bare metal}: The compiler still understands
            \bigbreak \noindent 
            \begin{cppcode}
                Widget* object = new Widget{};
            \end{cppcode}
            \bigbreak \noindent 
            Conceptually, this performs two operations:
            \begin{itemize}
                \item Obtain raw storage through operator new.
                \item Construct Widget in that storage.
            \end{itemize}
            \bigbreak \noindent 
            \begin{cppcode}
                void* memory = operator new(sizeof(Widget));
                Widget* object = ::new (memory) Widget{};
            \end{cppcode}
            \bigbreak \noindent 
            Whether ordinary new works depends on the runtime and linker configuration. In many bare-metal toolchains, the allocator ultimately obtains heap space through a low-level function such as \_sbrk, which you or the vendor runtime must provide.
            \bigbreak \noindent 
            You can also disable or forbid heap use by:
            \begin{itemize}
                \item Not implementing the required heap-growth routine
                \item Overriding global operator new
                \item Treating allocation calls as link errors
                \item Enforcing a project rule against dynamic containers
                \item Using linker symbols to reserve zero or limited heap space
            \end{itemize}
            ::new is a C++ expression that forces the compiler to call the global operator new function instead of any custom, class-specific operator new overrid
        \item \textbf{Placement new}: Placement new constructs an object in storage you already own:
            \bigbreak \noindent 
            \begin{cppcode}
                #include <new>
                #include <cstddef>

                alignas(Widget) std::byte storage[sizeof(Widget)];

                Widget* widget = new (storage) Widget{};

                widget->~Widget();
            \end{cppcode}
            \bigbreak \noindent 
            Placement new does not allocate memory. It only begins the object’s lifetime in the supplied storage.
            \bigbreak \noindent 
            Alignment and object lifetime must be handled correctly, so higher-level wrappers such as std::optional, std::variant, or a well-tested pool are generally preferable.



    \end{itemize}

    \pagebreak 
    \unsect{Bit manipulations}
    \begin{itemize}
        \item \textbf{Reverse bit pattern of a 32 bit integer (Iterative)}:  
            \bigbreak \noindent 
            \begin{cppcode}
                std::uint32_t rev(std::uint32_t n) {
                    std::uint32_t res{0}; 
                    for (size_t i{}; i < 32; ++i) {
                        res = res << 1 | (n & 1);
                        n >>= 1;
                    }

                    return res;
                }
            \end{cppcode}
        \item \textbf{Reverse bit pattern of a 32 bit integer (fast)}:
            \bigbreak \noindent 
            \begin{cppcode}
                unsigned int reverse_bits_fast(unsigned int n) {
                    // Swap adjacent single bits (01010101 and 10101010)
                    n = ((n & 0x55555555) << 1) | ((n & 0xAAAAAAAA) >> 1);
                    // Swap adjacent 2-bit pairs (00110011 and 11001100)
                    n = ((n & 0x33333333) << 2) | ((n & 0xCCCCCCCC) >> 2);
                    // Swap adjacent 4-bit nibbles (00001111 and 11110000)
                    n = ((n & 0x0F0F0F0F) << 4) | ((n & 0xF0F0F0F0) >> 4);
                    // Swap adjacent 8-bit bytes
                    n = ((n & 0x00FF00FF) << 8) | ((n & 0xFF00FF00) >> 8);
                    // Swap the 16-bit halves
                    n = (n << 16) | (n >> 16);
                    return n;
                }
            \end{cppcode}
        \item \textbf{Check if power of two}
            \bigbreak \noindent 
            \begin{cppcode}
                bool power_of_two(std::uint32_t value) {
                    if (value == 0) return false;
                    return value & (value-1) == 0;
                }
            \end{cppcode}
        \item \textbf{Clear lowest set bit}
            \bigbreak \noindent 
            \begin{cppcode}
                value & (value - 1)
            \end{cppcode}
        \item \textbf{Count number of set bits}
            \bigbreak \noindent 
            \begin{cppcode}
                unsigned count_set_bits(std::uint32_t value) {
                    unsigned count{};
                    while (value != 0) {
                        value &= value-1;
                        ++count;
                    }
                    return count;
                } 
            \end{cppcode}
        \item \textbf{Get msb}
            \bigbreak \noindent 
            \begin{cppcode}
                std::uint32_t get_msb(std::uint32_t value) {
                    value |= value >> 1;
                    value |= value >> 2;
                    value |= value >> 4;
                    value |= value >> 8;
                    value |= value >> 16;

                    return value ^= (value >> 1);
                }
            \end{cppcode}
        \item \textbf{Every value in a vector appears exactly twice except one, which appears once. Find that single value in O(n) time and O(1) extra space:}
            \bigbreak \noindent 
            \begin{cppcode}
            int find_single(const std::vector<int>& values);
            \end{cppcode}
            \bigbreak \noindent 
            The operator is bitwise XOR ($\land$). Two equal values cancel because $x \land x == 0$, and XOR with zero leaves a value unchanged. XOR is also associative and commutative, so the pairs need not be adjacent:
            \bigbreak \noindent 
            \begin{cppcode}
                int find_single(const std::vector<int>& values) {
                    int result = 0;
                    for (int ell : values) {
                        result ^= ell;
                    }

                    return result;
                }
            \end{cppcode}

    \end{itemize}

    \pagebreak 
    \unsect{Deducing Types}
    \begin{itemize}
        \item \textbf{Rule sets}: C++ has three distinct sets of type deduction rules
            \begin{enumerate}
                \item \textbf{Template type deduction}: Divided into three cases based on the parameter type
                    \begin{enumerate}
                        \item ParamType is a reference or pointer (but not a universal/forwarding reference).
                        \item ParamType is a universal/forwarding reference (T\&\&).
                        \item ParamType is passed by value (drops references, top-level const, and volatile)
                    \end{enumerate}
                \item \textbf{Auto}: It generally mirrors Template Type Deduction (the variable's type acts like ParamType and the initializer expression acts like the argument).
                    \bigbreak \noindent 
                    Braced initializer lists deduce to a std::initializer\_list under auto, which fails under standard template argument deduction
                \item \textbf{decltype / decltype(auto)}: Used to inspect the exact, unmodified type of a variable or expression.
                    \bigbreak \noindent 
                    Unlike auto and templates, decltype does not strip references or top-level cv-qualifiers; it gives you the exact declaration type or reference category (lvalue vs rvalue) of the expression
            \end{enumerate}
        \item \textbf{decltype(auto)}: decltype(auto) is a type specifier introduced in C++14 that tells the compiler to deduce a type using decltype rules rather than standard auto template deduction rules.
        \item \textbf{Template type deduction}: Consider
            \bigbreak \noindent 
            \begin{cppcode}
            template<typename T>
            void f(ParamType param);

            f(expr);
            \end{cppcode}
            \bigbreak \noindent 
            There are three cases for \texttt{ParamType}, each of which effect $T$s deduction differently
            \begin{enumerate}
                \item \textbf{ParamType} is a pointer or reference type. In this case, if the expr is a reference, the reference part is ignored. Then pattern-match expr’s type against ParamType to determine T.
                    \bigbreak \noindent 
                    \begin{cppcode}
                    template<typename T>
                    void f(T& param);

                    int x = 27;
                    const int cx = x;
                    const int& rx = x;
                    \end{cppcode}
                    \bigbreak \noindent 
                    Calling $f$ with $x,cx$, and $rx$ deduces $T$ as int, const int, const int.
                \item \textbf{ParamType} is a universal reference. If expr is an lvalue, both T and ParamType are deduced to be lvalue references. If expr is an rvalue, the “normal” (i.e., Case 1) rules apply.
                    \bigbreak \noindent 
                    Universal references are described heavily in a section above and will likely be described in more detail later.
                \item \textbf{ParamType} is neither a pointer nor a reference. As before, if expr’s type is a reference, ignore the reference part. If, after ignoring expr’s reference-ness, expr is const, ignore that, too. If it’s volatile, also ignore that. 
                    \bigbreak \noindent 
                    \begin{cppcode}
                        template<typename T>
                        void f(T param);

                        int x = 27;
                        const int cx = x;
                        const int& rx = x;
                    \end{cppcode}
                    \bigbreak \noindent 
                    Calling $f$ with $x,cx$, and $rx$ deduces $T$ as int in every call.
                    \bigbreak \noindent 
                    Note that even though cx and rx represent const values, param isn’t const. That makes sense. param is an object that’s completely independent of cx and rx—a copy of cx or rx. The fact that cx and rx can’t be modified says nothing about whether param can be. That’s why expr’s constness (and volatileness, if any) is ignored when deducing a type for param: just because expr can’t be modified doesn’t mean that a copy of it can’t be.
            \end{enumerate}
        \item \textbf{Array references}:
            \bigbreak \noindent 
            \begin{cppcode}
            int arr[10]; 
            int (&r)[10] = arr;
            \end{cppcode}
            \bigbreak \noindent 
            Here, $r$ is a reference to an array of 10 integers. With this, we can have access to the array in a function without being decayed to a pointer
            \bigbreak \noindent 
            \begin{cppcode}
                void f(int (&arr_ref)[10]) {
                    // 40
                    std::cout << sizeof(arr_ref) << "\n";
                }
            \end{cppcode}
            \bigbreak \noindent 
            Now, in
            \bigbreak \noindent 
            \begin{cppcode}
            template<typename T>
            void f(T& p);

            f(arr);
            \end{cppcode}
            \bigbreak \noindent 
            $T$ is deduced as int[10], and $p$ is int (\&)[10];
            \bigbreak \noindent 
            Interestingly, the ability to declare references to arrays enables creation of a template that deduces the number of elements that an array contains
            \bigbreak \noindent 
            \begin{cppcode}
                template<typename T, std::size_t N> constexpr std::size_t array_size(T (&)[N]) noexcept { return N; }
            \end{cppcode}
        \item \textbf{Displaying deduced}: Consider a class
            \bigbreak \noindent 
            \begin{cppcode}
            template<typename T>
            class TD;
            \end{cppcode}
            \bigbreak \noindent 
            Notice that TD does not use $T$, so it cannot be instantiated. Thus, the error message that is produced when we \textit{try} to instantiate with a type yields the type name that it tried to use. 
            \bigbreak \noindent 
            \begin{cppcode}
                std::vector<int> v{1,2,3};

                const int x = 20;
                TD<decltype(x)> t; // Error message discusses const int
                TD<decltype(v.begin()) w; // Error message discusses __gnu_cxx::__normal_iterator<int*, std::vector<int>>
            \end{cppcode}
        \item \textbf{Why prefer auto}: auto prevents uninitialized variables, avoids incorrect or duplicated type names, and adapts automatically when APIs change. It is especially useful for iterators, lambdas, and complicated template types.
            \bigbreak \noindent 
            Use an explicit type when it communicates important domain meaning that auto would obscure.
        \item \textbf{Forcing auto}: Some expressions return proxy types rather than ordinary values. For example, an element of std::vector<bool> may not be an actual bool.
            \bigbreak \noindent 
            Force the intended type when necessary
            \bigbreak \noindent 
            \begin{cppcode}
            auto enabled = static_cast<bool>(...)
            \end{cppcode}


    \end{itemize}

    \pagebreak 
    \unsect{Moving to modern C++}
    \begin{itemize}
        \item \textbf{Declare overriding functions with override}: override makes the compiler verify that a function really overrides a virtual base-class function. It detects mistakes involving const, reference qualifiers, parameter types, and exception specifications.
        \item \textbf{Prefer const\_iterator when modification is unnecessary}: Use cbegin, cend, and related operations when the iterator should not modify elements. This expresses intent and avoids unnecessarily granting write access.
        \item \textbf{Declare functions noexcept when they cannot throw}: noexcept is part of a function’s interface and can enable better optimization. Standard containers may move elements during reallocation only when their move operations are known not to throw.
            \bigbreak \noindent 
            Do not add it merely for performance: escaping exceptions from a noexcept function terminate the program.
        \item \textbf{Use constexpr whenever the operation legitimately supports compile-time use}: constexpr objects are constant and must be initialized by constant expressions. constexpr functions may execute during compilation or runtime, depending on their arguments and context.
        \item \textbf{Make const member functions thread-safe}: Multiple threads may call const operations on the same object. Therefore, a logically read-only operation should not perform unsynchronized writes to shared state.
            \bigbreak \noindent 
            Use a mutex for compound state and std::atomic for suitable independent values.
        \item \textbf{Understand automatic generation of special member functions}: The special member functions are:
            \begin{itemize}
                \item Default constructor.
                \item Destructor.
                \item Copy constructor.
                \item Copy assignment operator.
                \item Move constructor.
                \item Move assignment operator.
            \end{itemize}
            Declaring some of these can prevent others from being generated. In particular, user-declared copy operations or a destructor can suppress implicit move operations.
            \bigbreak \noindent 
            Prefer the Rule of Zero: let resource-owning member types manage copying, moving, and destruction.
        \item \textbf{Use std::unique\_ptr for exclusive ownership}: unique\_ptr is small, efficient, movable, and customizable with deleters. It should normally be the first ownership pointer considered.
            \bigbreak \noindent 
            It can also convert into a shared\_ptr if ownership later needs to become shared.
        \item \textbf{Use std::shared\_ptr for genuine shared ownership}: shared\_ptr maintains a reference-counted control block. The managed object is destroyed when the last owning pointer disappears.
            \bigbreak \noindent 
            Avoid constructing multiple independent shared\_ptr objects from the same raw pointer. Use std::enable\_shared\_from\_this when an owned object must safely produce another owner of itself.
        \item \textbf{Use std::weak\_ptr for non-owning observation}: A weak\_ptr observes an object managed by shared\_ptr without extending its lifetime. It can:
            \begin{itemize}
                \item Detect whether the object has expired.
                \item Temporarily acquire ownership using lock().
                \item Break cycles between shared\_ptr objects.
                \item Support caches and observer structures.
            \end{itemize}
        \item \textbf{Pimpl}: PImpl means Pointer to Implementation. A class stores a pointer to a private implementation object instead of storing its private data directly in the public class definition.
            It is primarily used to:
            \begin{itemize}
                \item Hide implementation details.
                \item Reduce header dependencies and rebuild times.
                \item Improve binary compatibility.
                \item Keep public headers small and stable.
            \end{itemize}
            Consider
            \bigbreak \noindent 
            \begin{cppcode}
                // widget.h
                #pragma once

                #include <string>
                #include <vector>

                class Widget {
                    public:
                    void add_name(std::string name);

                    private:
                    std::vector<std::string> names_;
                };
            \end{cppcode}
            \bigbreak \noindent 
            Every file including widget.h must parse the <vector> and <string> declarations. Changing the private representation can also cause every dependent source file to be recompiled.
            \bigbreak \noindent 
            Now, with Pimpl:
            \bigbreak \noindent 
            \begin{cppcode}
                // widget.h
                #pragma once

                #include <memory>
                #include <string>

                class Widget {
                public:
                    Widget();
                    ~Widget();

                    Widget(Widget&&) noexcept;
                    Widget& operator=(Widget&&) noexcept;

                    Widget(const Widget&);
                    Widget& operator=(const Widget&);
                    void add_name(std::string name);

                private:
                    class Impl;
                    std::unique_ptr<Impl> impl_;
                };
            \end{cppcode}
            \bigbreak \noindent 
            \begin{cppcode}
                // widget.cpp
                #include "widget.h"

                #include <utility>
                #include <vector>

                class Widget::Impl {
                public:
                    std::vector<std::string> names;
                };

                Widget::Widget()
                    : impl_(std::make_unique<Impl>()) { }

                Widget::~Widget() = default;

                Widget::Widget(Widget&&) noexcept = default;
                Widget& Widget::operator=(Widget&&) noexcept = default;

                Widget::Widget(const Widget& other)
                    : impl_(std::make_unique<Impl>(*other.impl_)) { }

                Widget& Widget::operator=(const Widget& other) {
                    if (this != &other) {
                        *impl_ = *other.impl_;
                    }

                    return *this;
                }

                void Widget::add_name(std::string name) {
                    impl_->names.push_back(std::move(name));
                }
            \end{cppcode}
            \bigbreak \noindent 
            Users of Widget know that Impl exists, but they do not know what it contains.
            \bigbreak \noindent 
            Notice that the destructor is defined in the .cpp file, at the end of widget.h, Impl is incomplete. std::unique\_ptr<Impl> can hold a pointer to an incomplete type. However, destroying that pointer eventually requires Impl to be complete. Therefore, this is normally incorrect:
            \bigbreak \noindent 
            \begin{cppcode}
                // widget.h
                ~Widget() = default;
            \end{cppcode}
            \bigbreak \noindent 
            The compiler may instantiate unique\_ptr<Impl>’s destructor where Impl is still incomplete. Instead, declare the destructor in the header, then define it after Impl is complete.
            \bigbreak \noindent 
            The same consideration often applies to move operations, so defining all special member functions in the .cpp file is a dependable convention.
            \bigbreak \noindent 
            Because std::unique\_ptr is noncopyable, a Pimpl class is noncopyable by default. We could either make move-only by deleting copy and copy assignment, or by implementing a deep copy.
        \item \textbf{Avoid overloading directly on forwarding references}: A forwarding-reference overload can match more arguments than expected and may defeat overloads accepting const T\&.
            \bigbreak \noindent 
            Forwarding-reference constructors are particularly dangerous because they can interfere with copying, moving, and derived-to-base conversions.
        \item \textbf{Notes about forwarding references}: 
            \bigbreak \noindent 
            \begin{cppcode}
            template<typename T>
            void f(T&& p);
            \end{cppcode}
            \bigbreak \noindent 
            It is attractive because it accepts both lvalues and rvalues and preserves their value category with std::forward
            \bigbreak \noindent 
            The problem is that this template accepts almost anything. It can capture arguments that should have gone to another overload, and invalid arguments often produce complicated errors deep inside the function.
            \bigbreak \noindent 
            \begin{cppcode}
                class Person {
                    public:
                        template<typename T>
                        explicit Person(T&& name)
                            : name_(std::forward<T>(name)) {}

                    private:
                        std::string name_;
                };
            \end{cppcode}
            \bigbreak \noindent 
            But copying a non-const Person may select the template instead of the copy constructor:
            \bigbreak \noindent 
            When the accepted type is known, write ordinary overloads:
            \bigbreak \noindent 
            \begin{cppcode}
            class Person {
            public:
                explicit Person(const std::string& name)
                    : name_(name) {}

                explicit Person(std::string&& name)
                    : name_(std::move(name)) {}

            private:
                std::string name_;
            };
            \end{cppcode}
            \bigbreak \noindent 
            The first overload copies from lvalues, the second moves from rvalues. This is explicit and prevents unrelated types from being accepted. The important distinction is that these are overloads for the concrete type std::string; replacing std::string with an unconstrained template parameter would reintroduce much of the original problem.
            \bigbreak \noindent 
            Often, a single by-value overload is simpler
            \bigbreak \noindent 
            \begin{cppcode}
            class Person {
            public:
                explicit Person(std::string name)
                    : name_(std::move(name)) {}

            private:
                std::string name_;
            };
            \end{cppcode}
            \bigbreak \noindent 
            For an lvalue, the argument is copied into name, and then moved into name\_. For an rvalue, the parameter is constructed or moved efficiently:
            \bigbreak \noindent 
            This is particularly useful when the function needs to own a copy of the argument anyway. The tradeoff is that an lvalue may involve one copy followed by one move. For types such as std::string, the additional move is normally inexpensive.
        \item \textbf{Tag dispatch}: Tag dispatch uses a compile-time property to choose an internal implementation.
            \bigbreak \noindent 
            Suppose integral arguments represent database IDs, while string-like arguments represent names:
            \bigbreak \noindent 
            \begin{cppcode}
            class Person {
            public:
                explicit Person(int id)
                    : Person(nameFromId(id)) {}

                template<typename T>
                explicit Person(T&& name)
                    : Person(
                          std::forward<T>(name),
                          std::is_integral<std::remove_reference_t<T>>{}
                      ) {}

            private:
                template<typename T>
                Person(T&& value, std::false_type)
                    : name_(std::forward<T>(value)) {}

                template<typename T>
                Person(T value, std::true_type)
                    : name_(nameFromId(value)) {}

                static std::string nameFromId(int id);

                std::string name_;
            };
            \end{cppcode}
            \bigbreak \noindent 
            The tags are
            \bigbreak \noindent 
            \begin{cppcode}
                std::true_type
                std::false_type
            \end{cppcode}
            \bigbreak \noindent 
            They select an overload at compile time.
            \bigbreak \noindent
            In C++17 and later, if constexpr is often clearer for simple cases:
            \bigbreak \noindent 
            \begin{cppcode}
                template<typename T>
                void process(T&& value) {
                    using U = std::remove_cvref_t<T>;

                    if constexpr (std::integral<U>) {
                        processId(value);
                    } else {
                        processName(std::forward<T>(value));
                    }
                }
            \end{cppcode}
        \item \textbf{Constrained templates}: C++20 concepts let you keep perfect forwarding while clearly limiting which types are accepted.
            \bigbreak \noindent 
            \begin{cppcode}
            #include <concepts>
            #include <string>
            #include <utility>

            class Person {
            public:
                template<typename T>
                    requires std::constructible_from<std::string, T>
                explicit Person(T&& name)
                    : name_(std::forward<T>(name)) {}

            private:
                std::string name_;
            };
            \end{cppcode}
            \bigbreak \noindent 
            A named concept can make the interface even clearer:
            \bigbreak \noindent 
            \begin{cppcode}
            template<typename T>
            concept StringLike =
                std::constructible_from<std::string, T>;

            class Person {
            public:
                template<StringLike T>
                explicit Person(T&& name)
                    : name_(std::forward<T>(name)) {}

            private:
                std::string name_;
            };
            \end{cppcode}






    \end{itemize}

    \pagebreak 
    \unsect{Modern STL}
    \begin{itemize}
        \item \textbf{std::ref}: In C++, std::ref is a helper function template that generates an object of type std::reference\_wrapper. It is used to pass true references to function templates (like std::thread or std::bind) or standard library algorithms that aggressively copy or consume their arguments by value      
            \bigbreak \noindent 
            Because native C++ references (T\&) are not objects, they cannot be copied, reassigned, or stored in standard containers directly. std::ref solves this by wrapping the reference inside a copyable, assignable object that behaves exactly like the underlying reference
            \bigbreak \noindent 
            \begin{cppcode}
                #include <functional>

                template<typename T>
                auto f(T p) -> void {
                    p.get() = 20;
                }

                int x = 10;
                f(std::ref(x));

                // 20
                std::cout << x << '\n';
            \end{cppcode}
        \item \textbf{std::cref}: std::cref stands for "constant reference". It is the exact same utility as std::ref, but instead of wrapping a mutable reference (T\&), it wraps a read-only, constant reference (const T\&).
            \bigbreak \noindent 
            You use std::cref when you want to pass an object to a value-copying template (like std::thread, std::bind, or std::async) to prevent copies and achieve better performance, but you also want a guarantee that the receiving function cannot modify your original variable.
            \bigbreak \noindent 
            If you attempt to modify a variable wrapped by std::cref, the compiler will block it, protecting your data:
        \item \textbf{std::bind and std::placeholders}: In C++, std::bind is a standard library function template defined in the <functional> header that generates a new callable object (a functor) by "binding" or fixing specific arguments of an existing callable (a function, lambda, or member function).
            \bigbreak \noindent 
            It allows for partial function application, where you pre-fill some function parameters now and leave others to be passed later at the call site
            \bigbreak \noindent 
            To leave arguments variable so they can be provided when the new function is actually called, you use placeholders (std::placeholders::\_1, std::placeholders::\_2, etc.)
            \bigbreak \noindent 
            \begin{cppcode}
                #include <functional>

                void printValues(int a, int b, int c) {
                    std::cout << "a: " << a << ", b: " << b << ", c: " << c << "\n";
                }

                using namespace std::placeholders;

                auto boundFunc = std::bind(printValues, 10, _1, _2);

                // Call site: 20 maps to _1, 30 maps to _2
                boundFunc(20, 30); // Output: a: 10, b: 20, c: 30
            \end{cppcode}
            \bigbreak \noindent 
            You can also reorder arguments using placeholders
            \bigbreak \noindent 
            \begin{cppcode}
                auto reorderedFunc = std::bind(printValues, _2, _1, 99);
                reorderedFunc(5, 15); // Output: a: 15, b: 5, c: 99
            \end{cppcode}
        \item \textbf{Binding member functions}: When binding a non-static class member function, you must pass a pointer to the member function as the first argument, and a pointer or reference to the object instance as the second argument
            \bigbreak \noindent 
            \begin{cppcode}
                struct Calculator {
                    void add(int x, int y) {
                        std::cout << "Sum: " << (x + y) << "\n";
                    }
                };

                Calculator calc;
    
                // Bind to a specific instance (&calc)
                auto quickAdd = std::bind(&Calculator::add, &calc, 5, std::placeholders::_1);
            \end{cppcode}
        \item \textbf{std::bind gotcha}: By default, std::bind copies or moves its arguments into the internal storage of the returned functional object. If your original function takes an argument by reference, std::bind will still copy it unless you explicitly wrap it using std::ref or std::cref
            \bigbreak \noindent 
            \begin{cppcode}
                void increment(int& num) { num++; }

                int main() {
                    int x = 0;

                    auto badIncrement = std::bind(increment, x); // Copies x
                    badIncrement(); // 'x' remains 0 inside main

                    auto goodIncrement = std::bind(increment, std::ref(x)); // Passes reference
                    goodIncrement(); // 'x' is now 1
                }
            \end{cppcode}
        \item \textbf{Modern std::bind functionality}: Prefer lambdas, they result in cleaner compilation errors, better optimization potential for compilers, and more readable code.
            \bigbreak \noindent 
            \begin{cppcode}
                void f(int x, int y, int z) {
                    std::cout << x << ", " << y << ", " << z << '\n';
                }

                auto g = [](int b, int c) { f(1,b,c); };
                g(2,3);
            \end{cppcode}



    \end{itemize}

    \pagebreak 
    \unsect{Exercises}
    \begin{itemize}
        \item \textbf{Fixed capacity ring buffer}: 
            \bigbreak \noindent 
            \begin{cppcode}
                class RingBuffer {
                    std::size_t capacity_{};
                    std::size_t size_{};

                    // Position of the next read
                    std::size_t read_{};

                    // Position of the next write
                    std::size_t write_{};

                    // Internal storage
                    std::vector<int> buf_{};
                public:
                    explicit RingBuffer(std::size_t capacity) : capacity_(capacity), buf_(capacity) {}; 

                    bool push(int);
                    std::optional<int> pop();
                    std::size_t size() const noexcept;
                    bool full() const noexcept;
                    void print() const;
                };
            \end{cppcode}
            \bigbreak \noindent 
            \begin{cppcode}
                bool RingBuffer::push(int value) {
                    if (full()) return false;

                    buf_[write_] = value;
                    write_ = (write_ + 1) % capacity_;
                    ++size_;

                    return true;
                }
            \end{cppcode}
            \bigbreak \noindent 
            \begin{cppcode}
                std::optional<int> RingBuffer::pop() {
                    if (size() == 0) return std::nullopt;

                    int head = buf_[read_];
                    read_ = (read_ + 1) % capacity_;
                    --size_;

                    return head;
                }
            \end{cppcode}
            \bigbreak \noindent 
            \begin{cppcode}
                std::size_t RingBuffer::size() const noexcept {
                    return size_;
                }

                bool RingBuffer::full() const noexcept {
                    return size_ == capacity_;
                }
            \end{cppcode}
            \bigbreak \noindent 
            \begin{cppcode}
                void RingBuffer::print() const {
                    std::size_t idx = read_;

                    for (size_t i{0}; i < size_; ++i) {
                        if (i != 0) std::cout << ", ";
                        std::cout << buf_[idx];
                        idx = (idx + 1) % capacity_;
                    }

                    std::cout << "\n";
                }
            \end{cppcode}

        \item \textbf{Merge intervals}: Implement a function that merges overlapping closed intervals:
            \bigbreak \noindent 
            \begin{cppcode}
                struct Interval {
                    int start;
                    int end;
                };

                std::vector<Interval> merge_intervals(
                const std::vector<Interval>& intervals
                );
            \end{cppcode}
            \bigbreak \noindent 
            Each input interval satisfies start <= end. The input may be unsorted. Return the merged intervals sorted by start, without modifying the input.
            \bigbreak \noindent 
            \begin{cppcode}
                Input:  [5, 7], [1, 3], [3, 4], [10, 12]
                Output: [1, 4], [5, 7], [10, 12]
            \end{cppcode}
            \bigbreak \noindent 
            \begin{cppcode}
            std::vector<Interval> merge_intervals(const std::vector<Interval>& intervals) {
                if (intervals.empty()) return {};

                std::vector<Interval> ret;
                std::vector<Interval> cpy = intervals;

                if (cpy.size() == 1) return cpy;

                std::sort(cpy.begin(), cpy.end(), [](const Interval& a, const Interval& b) {
                    return a.start < b.start;
                });

                std::size_t left = 0;
                while (left < cpy.size()) {
                    Interval merged = cpy[left];

                    std::size_t last_merged = left;
                    for (std::size_t i{left+1}; i < cpy.size(); ++i) {
                        if (cpy[i].start <= merged.end) {
                            merged.end = std::max(merged.end, cpy[i].end);
                            last_merged = i;
                            continue;
                        }

                        break;
                    }

                    if (last_merged == left) ++left;
                    else left = last_merged+1;

                    ret.push_back(merged);
                }

                return ret;
            }
            \end{cppcode}


    \end{itemize}
   

\end{document}
